% Twin correlations of truncated von Mangoldt functions up to level 9/17:
% a formalized proof and the barrier at 4/7.
%
% DRAFT.  This manuscript was prepared with extensive use of generative AI;
% see the section "AI usage disclosure" at the end and AI_DISCLOSURE.md in
% the repository root.  It has NOT yet been verified by a human expert.
%
% Add affiliation and contact details before any circulation.
\documentclass[11pt]{amsart}

\usepackage[T1]{fontenc}
\usepackage[utf8]{inputenc}
\usepackage{amsmath,amssymb,amsthm,mathtools}
\usepackage{microtype}
\usepackage[margin=1.15in]{geometry}
\usepackage[colorlinks=true,citecolor=blue,linkcolor=blue,urlcolor=blue]{hyperref}

\newtheorem{theorem}{Theorem}[section]
\newtheorem{proposition}[theorem]{Proposition}
\newtheorem{lemma}[theorem]{Lemma}
\newtheorem{corollary}[theorem]{Corollary}
\theoremstyle{definition}
\newtheorem{definition}[theorem]{Definition}
\theoremstyle{remark}
\newtheorem{remark}[theorem]{Remark}
\newtheorem{claim}[theorem]{Claim}

\numberwithin{equation}{section}

\newcommand{\N}{\mathbb N}
\newcommand{\Z}{\mathbb Z}
\newcommand{\Q}{\mathbb Q}
\newcommand{\R}{\mathbb R}
\newcommand{\C}{\mathbb C}
\newcommand{\e}{\mathrm e}
\newcommand{\cS}{\mathcal S}
\newcommand{\cF}{\mathcal F}
\newcommand{\cW}{\mathcal W}
\newcommand{\cX}{\mathcal X}
\newcommand{\cB}{\mathcal B}
\newcommand{\eps}{\varepsilon}
\newcommand{\inv}[1]{\overline{#1}}
\newcommand{\sumstar}{\sideset{}{^*}\sum}

\title[Twin correlations of $\Lambda_R$: a formalized proof and the $4/7$ barrier]
{Twin correlations of truncated von Mangoldt functions\\ up to level $9/17$:\\
a formalized proof and the barrier at $4/7$}

\author[Zach Conn]{Investigator: Zach Conn}
\thanks{Zach Conn is the investigator responsible for this work.  The
mathematical content was generated with extensive assistance from
generative AI systems; see the AI usage disclosure at the end of the
paper.  This is a draft that has not yet been checked by a human expert.}

\date{Draft of 15 November 2026}

\subjclass[2020]{Primary 11N37; Secondary 11L05, 11L07, 11B30}
\keywords{M\"obius function, Kloosterman fractions, determinant equation,
Heath-Brown identity, additive energy of rationals, correlations of divisor
sums, formal verification, Lean, additive divisor problem}

\begin{document}

\begin{abstract}
Let $\Lambda_R(n)=\sum_{d\mid n,\,d\le R}\mu(d)\log(R/d)$ be the truncated
von Mangoldt function, so that $\Lambda_N(n)=\Lambda(n)$ for $1<n\le N$.  We
give a proof, formalized in Lean~4 with Mathlib, that
\[
 \sum_{n\le N}\Lambda_R(n)\Lambda_R(n+2)
 =\mathfrak S\,N+O\bigl(N\exp(-c\sqrt{\log R})\bigr)
\]
uniformly for $2\le R\le N^{\theta}$, for every fixed $\theta<9/17$.  Here
$\mathfrak S$ is the twin-prime constant.  The formal proof has no hypotheses
and depends only on Lean's standard axioms.  It includes formal proofs of
the Bloom--Maynard bound for the additive energy of rationals and of the
prime number theorem for $\mu$ with classical error term.

The exponent $9/17$ is not the best that existing methods give.  The
estimates of Deshouillers and Iwaniec, used as in Conrey's work on
mollifiers, appear to reach $\theta<4/7$.  Our proof is of a different kind.
It uses no spectral theory: its inputs are Weil's bound, the Heath-Brown
identity, a Korolev--Changa bilinear inequality averaged over moduli, and the
Bloom--Maynard bound.  This is what made a complete formalization feasible.

The second part of the paper is a synthesis of what is known for this
problem past level $1/2$, and a report on an attempt to pass $4/7$.
\begin{itemize}
\item We reduce the problem of passing $4/7$ to one estimate, a restricted
additive divisor problem for $d_3\times d_3$.
\item We give equivalent forms of that estimate: a mean square of
exponential sums at a single level, a second moment of automorphic
$L$-functions twisted by an Atkin--Lehner sign, and Weyl sums over roots of
a quadratic congruence.
\item We show that three different methods are tight at exactly $4/7$, and
report numerical experiments.
\end{itemize}
We do not pass $4/7$.  This manuscript was prepared with substantial AI
assistance and has not yet been verified by a human expert; see the
disclosure at the end.
\end{abstract}

\maketitle

\section{Introduction}\label{sec:intro}

\subsection{Correlations of truncated von Mangoldt functions}
The twin-prime conjecture, in its quantitative Hardy--Littlewood form,
asserts that
\[
 \sum_{n\le N}\Lambda(n)\Lambda(n+2)\sim\mathfrak S N,
 \qquad
 \mathfrak S=2\prod_{p>2}\Bigl(1-\frac1{(p-1)^2}\Bigr)=1.3203\ldots.
\]
Since $\Lambda(n)=\sum_{d\mid n}\mu(d)\log(n/d)$, a natural model for
$\Lambda$ is the truncated divisor sum
\[
 \Lambda_R(n)=\sum_{\substack{d\mid n\\ d\le R}}\mu(d)\log(R/d),
\]
which satisfies $\Lambda_N(n)=\Lambda(n)$ for $1<n\le N$.  The truncation
level $R=N^\theta$ therefore acts as a dial on the correlation
\[
 \mathcal C(N,R)=\sum_{n\le N}\Lambda_R(n)\Lambda_R(n+2):
\]
at $\theta=1$ it is the twin-prime sum itself.  For $\theta\le1/2$, every
pair of divisors $d,e\le R$ involved has $de\le N$.  The Chinese remainder
theorem then gives $\mathcal C(N,R)\sim\mathfrak SN$; see, e.g.,~\cite{GY}.
Truncated correlations of this kind are basic objects in the work of
Goldston and Y\i ld\i r\i m and in the Goldston--Pintz--Y\i ld\i r\i m
method~\cite{GY,GPY}.  For the Selberg sieve weights $\lambda_z$, Walker
proved a fixed-shift asymptotic for $\sum_n\lambda_z(n)\lambda_z(n+k)$ that is
nontrivial for $z\le X^{27/53-\eps}$~\cite{Walker}.  He obtained it from the
Bettin--Chandee bound with the frequency held fixed.  He then applied it to
conditional results on the variance of primes in short intervals and on the
pair correlation of zeros of the Riemann zeta function.

Past $\theta=1/2$ one meets pairs with $de>N$, and the problem becomes one
about sums of Kloosterman fractions.  It is a binary analogue of passing
level $1/2$ in theorems of Bombieri--Vinogradov type~\cite{BFI}.
\begin{itemize}
\item Coppola and Laporta~\cite{CL} treat the fixed-shift autocorrelation of
an arbitrary sieve function of level $\theta$, using the bilinear
Kloosterman-fraction estimate of Duke, Friedlander and Iwaniec~\cite{DFI}.
They obtain an asymptotic formula for $\theta<48/95=1/2+1/190$; $\Lambda_R$
is a special case.
\item Inserting the stronger trilinear estimate of Bettin and
Chandee~\cite{BC} into the same reduction gives
$\theta<17/33=1/2+1/66$ in the critical range (Appendix~\ref{app:BC}).
\end{itemize}
Both arguments treat the coefficients $\mu(d)\log(R/d)$ as an arbitrary
bounded sequence.  Our main result exploits their arithmetic structure.

The same sums of Kloosterman fractions occur in the off-diagonal terms of the
twisted second moment of the Riemann zeta function with a mollifier of
length $T^\theta$.  There the coefficients have the shape
$\mu(d)P(\log(R/d)/\log R)$.
\begin{itemize}
\item Balasubramanian, Conrey and Heath-Brown~\cite{BCHB} reached
$\theta<9/17$ for mollifier coefficients.
\item Conrey~\cite{Con89} reached $\theta<4/7$, using the Heath-Brown
identity for $\mu$, Weil's bound and the estimates of Deshouillers and
Iwaniec~\cite{DI82}; see also~\cite{PR18}.
\item For arbitrary coefficients, Bettin, Chandee and
Radziwi\l\l~\cite{BCR} obtain $17/33$ from~\cite{BC} and $48/95$
from~\cite{DFI}, the same pair of exponents as above.
\end{itemize}
In a related direction, Fouvry and Shparlinski~\cite{FS11} and
Irving~\cite{Irv14} bound Kloosterman sums over primes on average over the
modulus, using a combinatorial decomposition of $\Lambda$.  See
Remark~\ref{rem:comparison}.  Fouvry and Radziwi\l\l~\cite{FR22} prove a
level of distribution $17/33$ for unbalanced convolutions, which covers
sieve weights to very large moduli; it does not apply directly to a fixed
shift.

\begin{theorem}\label{thm:LambdaR}
Let $0<\theta<9/17$.  There are an absolute constant $c>0$ and a constant
$C_\theta$ such that, for all $N\ge2$ and $2\le R\le N^\theta$,
\[
 \Bigl|\sum_{n\le N}\Lambda_R(n)\Lambda_R(n+2)-\mathfrak S N\Bigr|
 \le C_\theta\,N\exp\bigl(-c\sqrt{\log R}\bigr).
\]
\end{theorem}

\emph{Status of the exponent.}  The exponents recorded explicitly for this
fixed-shift problem are $48/95=1/2+1/190$ for $\Lambda_R$~\cite{CL} and
$27/53$ for Selberg weights~\cite{Walker}.  Theorem~\ref{thm:LambdaR}
exceeds both, but $9/17$ is not the best exponent that existing methods
give:
\begin{itemize}
\item Irving's estimate~\cite[Theorem~1.3]{Irv14}, transposed from $\Lambda$
to $\mu$, appears to give Theorems~\ref{thm:LambdaR} and~\ref{thm:level} for
$\theta<8/15$ (Remark~\ref{rem:comparison});
\item Conrey's method~\cite{Con89}, through the estimates of Deshouillers
and Iwaniec~\cite{DI82}, appears to give them for $\theta<4/7$
(Claim~\ref{claim:conrey}).
\end{itemize}
Neither transposition has been checked by a human, and neither is
formalized.  We therefore do not present $9/17$ as a record.

\subsection{What this paper contributes}\label{sec:contrib}
\begin{itemize}
\item \emph{A formalized proof.}  Theorems~\ref{thm:LambdaR},
\ref{thm:level} and~\ref{thm:main} are proved in Lean~4 with no hypotheses
(Section~\ref{sec:lean}).  To our knowledge this is the first machine-checked
asymptotic for a correlation of truncated divisor sums past the level $1/2$
at which the Chinese remainder theorem stops.  The proof was chosen for
this: it avoids the spectral theory of automorphic forms, none of which is
available in a proof assistant at present.
\item \emph{A different mechanism.}  The proof joins the Poisson frequency
to the shorter factor of $\mu$ and controls the resulting fractions by the
Bloom--Maynard energy bound, after averaging a Korolev--Changa inequality
over the modulus (Theorem~\ref{thm:meanvalue}).  It is weaker than the
spectral method on the diagonal and incomparable with Irving's estimate off
it (Remark~\ref{rem:comparison}).
\item \emph{A synthesis of the problem past $1/2$.}
Section~\ref{sec:barrier} collects what the literature gives, isolates the
one estimate whose absence stops every method at $4/7$, gives equivalent
forms of it, and reports numerical experiments and the estimates we examined
that do not apply.  That section contains no new theorem about level $4/7$.
\end{itemize}

\subsection{M\"obius-weighted divisor pairs}
For $d,e\ge1$ let
\[
 \varrho(d,e)=\begin{cases}1/[d,e],&(d,e)\mid 2,\\0,&\text{otherwise},\end{cases}
 \qquad
 r_N(d,e)=\#\{1\le n\le N:\ d\mid n,\ e\mid n+2\}-N\varrho(d,e).
\]
Here $\varrho(d,e)$ is the density of the integers $n$ with $d\mid n$ and
$e\mid n+2$, and $|r_N(d,e)|\le1$.  Theorem~\ref{thm:LambdaR} is deduced
in Section~\ref{sec:LambdaR} from the following estimate together with the
prime number theorem for $\mu$.

\begin{theorem}\label{thm:level}
For every $\theta<9/17$ there is $\delta=\delta(\theta)>0$ such that, for all
$N\ge2$ and all integers $1\le K,L\le N^\theta$,
\[
 \sum_{d\le K}\sum_{e\le L}\mu(d)\mu(e)\,r_N(d,e)\ll_\theta N^{1-\delta}.
\]
\end{theorem}

The trivial bound is $KL$, so the content of Theorem~\ref{thm:level} lies
in the range $KL>N$.  There it is a level of distribution beyond $1$ for the
products $de$, which may be as large as $N^{18/17-\eps}$.

\subsection{The determinant sum}
Theorem~\ref{thm:level} rests on an estimate for the determinant sums that
arise when $n=ad$ and $n+2=be$.  Let $W\colon\R^4\to\C$ be a fixed smooth
function supported in $[1,2]^4$.  Let $X\ge 3$ and $A,B\ge 1$, put $D=X/A$
and $E=X/B$, and let $\kappa\in\{2,-2\}$.  Define
\begin{equation}\label{eq:S}
 \cS_\kappa=\cS_\kappa(W;A,B,D,E)
 =\sum_{\substack{a,b,d,e\ \mathrm{odd}\\ be-ad=\kappa}}
 \mu(d)\mu(e)\,
 W\!\left(\frac aA,\frac bB,\frac dD,\frac eE\right).
\end{equation}
Since $a,b,d,e$ are odd, every common divisor of $d$ and $e$ divides
$\kappa$, so $(d,e)=1$.  (For the prime analogue $p_1d-p_2b=r$ see
Guria~\cite{Gur}.)  Write
\begin{equation}\label{eq:xy}
 A=X^{x},\qquad B=X^{y},\qquad m=\min(x,y),\qquad M=\max(x,y).
\end{equation}
The ``hard square'' is $0<x,y\le 1/2$, in which $d$ and $e$ both have size
at least $\sqrt X$.  Put
\[
 \Omega=\bigl\{(x,y)\in(0,1/2]^2:\ 8m+9M>8\bigr\},
 \qquad
 \sigma(x,y)=\frac{8m+9M-8}{12}.
\]
Since $M\le1/2$, every point of $\Omega$ has $m>7/16$.

\begin{theorem}\label{thm:main}
Let $K\subset\Omega$ be compact, let $\kappa\in\{2,-2\}$ and let $\eps>0$.
There is an explicit term $\cS_\kappa^{0}$, given in
\eqref{eq:zero-mode} below, such that, uniformly for $(x,y)\in K$,
\begin{equation}\label{eq:main-power}
 \cS_\kappa=\cS_\kappa^{0}
 +O_{K,W,\eps}\bigl(X^{1-\sigma(x,y)+\eps}\bigr).
\end{equation}
Moreover there is an absolute constant $c>0$ such that
$\cS_\kappa^{0}\ll_{W}X\exp(-c\sqrt{\log X})$.  Consequently
\begin{equation}\label{eq:main-log}
 \cS_\kappa\ll_{K,W}X\exp\bigl(-c\sqrt{\log X}\bigr).
\end{equation}
\end{theorem}

The term $\cS_\kappa^0$ is the zero frequency after Poisson summation, a
smooth double sum of $\mu(d)\mu(e)/e$.  The power saving
\eqref{eq:main-power} concerns everything else.  In the correlation problem
of Theorem~\ref{thm:level}, the zero frequency is exactly the expected count
$N\varrho(d,e)$ and is subtracted; this is why Theorem~\ref{thm:level} has a
power saving.

\begin{corollary}\label{cor:diagonal}
If $x=y$, the conclusion \eqref{eq:main-log} holds whenever $x>8/17$, that
is, whenever $D=E\le X^{9/17-\eps}$ for some fixed $\eps>0$.
\end{corollary}

\begin{remark}[Comparison]\label{rem:comparison}
In Appendix~\ref{app:BC} we insert the unconditional trilinear bound of
Bettin and Chandee~\cite{BC} into the reduction of Section~\ref{sec:proof},
in place of Theorem~\ref{thm:meanvalue}.  This gives cancellation in the
region
\[
 16m+11M>13,\qquad 13m+20M>16 .
\]
On the diagonal this is $x>16/33$, that is, $D=E\le X^{17/33-\eps}$.  The
region $\Omega$ strictly contains it, because
\[
 8m+9M=\tfrac{43}{177}(16m+11M)+\tfrac{56}{177}(13m+20M)
\]
and $\tfrac{43}{177}\cdot 13+\tfrac{56}{177}\cdot16=\tfrac{1455}{177}>8$.
The boundary of $\Omega$ inside the ordered half-square is the segment from
$(7/16,1/2)$ to $(8/17,8/17)$.

Irving~\cite[Theorem~1.3]{Irv14} bounds
$\sum_{q\sim Q}|\sum_{p\sim x}e(a\bar p/q)|$ for a fixed frequency $a$.  His
proof uses no property of the primes beyond Vaughan's identity, which has
the same Type~I/II shape for $\mu$.  Transposed to $\mu$, and inserted
frequency by frequency in place of Theorem~\ref{thm:meanvalue}, it bounds
the nonzero frequencies by $X^{\eps}(E^{1/2}D^{11/8}+E^{7/6}D^{2/3})$.  This
gives cancellation when
\[
 11m+4M>7,\qquad 4m+7M>5 .
\]
On the diagonal this is $x>7/15$, which is better than $\Omega$; at
$M=1/2$ it is $m>5/11$, which is worse than $m>7/16$.  The two boundaries
cross at $(31/67,32/67)$.  In the correlation problem of
Theorem~\ref{thm:level} the same input would give $\theta<8/15$.  This
transposition is an exponent computation that has not been checked in
detail; see Section~\ref{sec:barrier} for its place among the other
exponents.
\end{remark}

\subsection{What is and is not new}\label{sec:significance}
The M\"obius function is not an arbitrary sequence, and its structure can be
used to pass level $1/2$, as in the mollifier
problem~\cite{BCHB,Con89}.  Our mechanism has three steps:
\begin{itemize}
\item factor $\mu(d)$ by the Heath-Brown identity;
\item join the Poisson frequency to the shorter factor;
\item control the resulting fractions by the Bloom--Maynard energy bound,
after averaging a Korolev--Changa inequality over the modulus.
\end{itemize}
Combining a combinatorial decomposition of $\mu$ with Kloosterman-fraction
estimates averaged over the modulus is classical~\cite{Con89,FS11,Irv14,PR18}.
What appears to be new in the proof is the use of the Bloom--Maynard energy
bound for the fractions $h/u$, together with a two-sided Korolev--Changa
inequality averaged over the modulus.  The mechanism is isolated in a
mean-value theorem, Theorem~\ref{thm:meanvalue}, which is uniform in the
frequency coefficients.

The limits should be stated equally clearly.
\begin{itemize}
\item The exponent is not a record; see the status paragraph after
Theorem~\ref{thm:LambdaR}.
\item Methods of this kind are subject to the parity problem.  The dial
moves from $0.5$ to about $0.53$, or to about $0.57$ with the spectral
method, and the twin-prime case $\theta=1$ is no closer in any qualitative
sense.
\item A literature check found no earlier fixed-shift result stated for
$\Lambda_R$ past $48/95$; the results we know of~\cite{CL,Walker} use
general-coefficient bounds.  The better exponents come from transposing
methods built for other problems~\cite{BCHB,Con89,Irv14,PR18}.
\end{itemize}

\subsection{Method}
Suppose $x\le y$, so that $D\ge E$.  Poisson summation in $a$ modulo $2e$
turns $\cS_\kappa$ into a zero-frequency term plus a trilinear form
\[
 \frac AE\sum_{e\asymp E}\mu(e)\sum_{d\asymp D}\mu(d)
 \sum_{0<|h|\lesssim T}\nu_h\, e_e(h\inv d),
 \qquad T=\frac{E}{A}=X^{1-x-y}.
\]
A single application of Cauchy--Schwarz over $e$ removes $\mu(e)$.  This
reduces Theorem~\ref{thm:main} to a mean-value estimate over moduli,
Theorem~\ref{thm:meanvalue}, for
\[
 \cX=\Bigl(\sum_{E/2<e\le2E}\Bigl|\sum_{(d,2e)=1}\mu(d)V(d/D)
 \sum_{|h|\le H}\nu_h\,e_e(h\inv d)\Bigr|^2\Bigr)^{1/2}.
\]
We decompose $\mu(d)$ by the Heath-Brown identity.  For balanced
factorizations $d=uw$ with $D^{1/3}\ll u\ll D^{1/2}$, we join the frequency
$h$ to the shorter factor $u$.  This produces a family of reduced fractions
$h/u$ and the phase $e_q(\inv z\,\inv w)$ with $\inv z=h\inv u$.  We then
apply a Korolev--Changa bilinear inequality~\cite{KC} \emph{before} the
reciprocal energies are bounded pointwise, and average those energies over
the modulus.  The energy of the integers $w$ is handled by divisor counting.
The energy of the fractions $h/u$ is controlled by the theorem of Bloom and
Maynard~\cite{BM} on additive energy of rationals with bounded
denominators.  Boxes with a long smooth variable are handled by Poisson
summation and Weil's bound.  The constraint $8m+9M>8$ is exactly the
condition $T<D^{1/9}$, under which the balanced estimate beats the trivial
scale.

The same mean-value theorem gives Theorem~\ref{thm:level}.  Poisson
summation in $n$ modulo $[d,e]$ produces the same Kloosterman fractions
directly (Section~\ref{sec:level}).  Theorem~\ref{thm:LambdaR} then follows
by Abel summation and an evaluation of the main term
(Section~\ref{sec:LambdaR}).

\subsection{Formal verification}
Theorems~\ref{thm:LambdaR} and~\ref{thm:main} and
Corollary~\ref{cor:diagonal} have been formalized in Lean~4 with Mathlib,
with no hypotheses; see Section~\ref{sec:lean}.  The formal proofs cover:
\begin{itemize}
\item the determinant sums (Theorem~\ref{thm:main}), with the exact saving
$\sigma(x,y)$;
\item the level of distribution (Theorem~\ref{thm:level});
\item the mean-value theorem (Theorem~\ref{thm:meanvalue});
\item the Bloom--Maynard energy bound;
\item the prime number theorem for $\mu$ with classical error term;
\item the deduction of Section~\ref{sec:LambdaR}.
\end{itemize}
It reuses existing formalizations of the Weil bound for Kloosterman sums,
the divisor bound, and the zero-free region of~$\zeta$.  The exponent
arithmetic and the elementary identities of Section~\ref{sec:barrier} are
also machine-checked; the analytic estimates quoted there are not.

\subsection{Scope}
We make no claim about the twin-prime conjecture itself.  We do not claim
that $9/17$ is the best exponent available, and we do not pass $4/7$.

\section{Notation and standard lemmas}\label{sec:prelim}

We write $e(t)=\e^{2\pi i t}$ and $e_q(t)=e(t/q)$.  For an integer $x$
coprime to $q$, $\inv x$ denotes an inverse of $x$ modulo $q$; the modulus
is always clear from context.  We write $\tau$ for the divisor function and
$\tau_k$ for the $k$-fold divisor function.  We also use $(m,n)$ for the
greatest common divisor and $\|\cdot\|_p$ for $\ell^p$ norms of finitely
supported sequences.  A smooth dyadic partition of unity is a smooth
$\rho_0$ supported in $[1/2,2]$ with $\sum_{k\in\Z}\rho_0(t/2^k)=1$ for all
$t>0$.  For a smooth function $F$ on a compact set, $\|F\|_{C^k}$ is the
maximum of the sup norms of its partial derivatives of order at most $k$.

\begin{lemma}[Divisor bounds]\label{lem:divisor}
For every $\eps>0$, $\tau(n)\ll_\eps n^\eps$.  For fixed $k\ge1$,
$\sum_{n\le N}\tau_k(n)^2\ll_k N(\log 2N)^{k^2-1}$.
\end{lemma}

\begin{lemma}[Poisson summation in a residue class]\label{lem:poisson}
Let $f\in C_c^\infty(\R)$, $q\ge1$ and $c\in\Z$.  Then
\[
 \sum_{n\equiv c\ (\mathrm{mod}\ q)}f(n)
 =\frac1q\sum_{h\in\Z}\widehat f\Bigl(\frac hq\Bigr)e\Bigl(\frac{hc}{q}\Bigr),
 \qquad \widehat f(\xi)=\int_\R f(t)e(-\xi t)\,dt .
\]
\end{lemma}

\begin{lemma}[Weil--Estermann]\label{lem:weil}
For $q\ge1$ and $m,n\in\Z$, the Kloosterman sum
$S(m,n;q)=\sumstar_{x\bmod q}e_q(mx+n\inv x)$ satisfies
$|S(m,n;q)|\le\tau(q)(m,n,q)^{1/2}q^{1/2}$.
\end{lemma}

This is Weil's bound~\cite{Weil} for prime moduli, extended to all moduli
by Estermann~\cite{Est}; see also~\cite[Ch.~11]{IK}.

\begin{lemma}[Incomplete Kloosterman sums]\label{lem:incomplete}
Let $q\ge1$, let $(b,q)=1$, and let $f$ be smooth and supported in
$[N,cN]$ with $|f^{(j)}|\le L N^{-j}$ for $0\le j\le2$.  Then
\[
 \sum_{(v,q)=1}f(v)\,e_q(b\inv v)\ll_c L\,\tau(q)\bigl(q^{1/2}+Nq^{-1/2}\bigr).
\]
\end{lemma}

\begin{proof}
Split $v$ into reduced residue classes $r$ modulo $q$ and apply
Lemma~\ref{lem:poisson} in each class.  This gives
$q^{-1}\sum_{\ell}\widehat f(\ell/q)S(\ell,b;q)$.  Since $(b,q)=1$,
Lemma~\ref{lem:weil} gives $|S(\ell,b;q)|\le\tau(q)q^{1/2}$.  Two
integrations by parts give
$|\widehat f(\xi)|\ll LN\min(1,(N|\xi|)^{-2})$, so
$\sum_\ell|\widehat f(\ell/q)|\ll LN(1+q/N)$.
\end{proof}

\begin{lemma}[M\"obius sums with a coprimality condition]\label{lem:mobius}
There is an absolute constant $c>0$ such that, uniformly for $x\ge2$ and
integers $1\le q\le x^{100}$,
\[
 M_q(x):=\sum_{\substack{n\le x\\(n,q)=1}}\mu(n)\ll x\exp\bigl(-c\sqrt{\log x}\bigr).
\]
Consequently, if $G$ is smooth and supported in a fixed compact subset of
$(0,\infty)$, then uniformly for $1\le q\le N^{100}$
\[
 \sum_{(n,q)=1}\mu(n)G(n/N)\ll_G N\exp\bigl(-c\sqrt{\log N}\bigr).
\]
\end{lemma}

\begin{proof}
Comparing Euler products,
$\sum_{(n,q)=1}\mu(n)n^{-s}=\zeta(s)^{-1}\sum_{k\mid q^\infty}k^{-s}$, where
$k\mid q^\infty$ means that every prime factor of $k$ divides $q$.  Hence
$M_q(x)=\sum_{k\mid q^\infty,\,k\le x}M(x/k)$, where
$M(y)=\sum_{n\le y}\mu(n)$.  The classical prime number theorem for the
M\"obius function gives $M(y)\ll y\exp(-c_0\sqrt{\log y})$ for $y\ge2$;
see~\cite[Ch.~6]{MV}.  Also $|M(y)|\le y$ trivially.

For $k\le\sqrt x$ we get a contribution
$\ll x\exp(-c_0\sqrt{\tfrac12\log x})\sum_{k\mid q^\infty}k^{-1}$.  The sum
equals $q/\varphi(q)\ll\log\log(q+2)$.  For $k>\sqrt x$ we use the trivial
bound and Rankin's trick:
\[
 \sum_{\substack{k\mid q^\infty\\k>\sqrt x}}\frac xk
 \le x^{3/4}\sum_{k\mid q^\infty}k^{-1/2}
 =x^{3/4}\prod_{p\mid q}(1-p^{-1/2})^{-1}
 \le x^{3/4}\exp\Bigl(4\sum_{p\mid q}p^{-1/2}\Bigr).
\]
Here $\sum_{p\mid q}p^{-1/2}\ll\sqrt{\omega(q)}\ll\sqrt{\log x}$, so this
part is $\ll x^{3/4}\exp(O(\sqrt{\log x}))$.  The first claim follows.  The
second follows by partial summation:
$\sum_{(n,q)=1}\mu(n)G(n/N)=-\int M_q(t)G'(t/N)\,dt/N$, and the integrand
is supported on $t\asymp N$.
\end{proof}

\begin{lemma}[Heath-Brown identity]\label{lem:HB}
Let $J\ge1$, $Z\ge1$, and $1\le n\le Z^J$.  Then
\[
 \mu(n)=\sum_{j=1}^{J}(-1)^{j-1}\binom Jj
 \sum_{\substack{m_1\cdots m_j n_1\cdots n_{j-1}=n\\ m_1,\dots,m_j\le Z}}
 \mu(m_1)\cdots\mu(m_j).
\]
\end{lemma}

\begin{proof}
Let $M_Z(s)=\sum_{m\le Z}\mu(m)m^{-s}$.  The Dirichlet coefficients of
$1-\zeta(s)M_Z(s)$ vanish for $n\le Z$, so those of $(1-\zeta M_Z)^J$
vanish for $n\le Z^J$, and so do those of $\zeta^{-1}(1-\zeta M_Z)^J$.
Expand by the binomial theorem:
\[
 \zeta^{-1}(1-\zeta M_Z)^J
 =\zeta^{-1}+\sum_{j\ge1}\binom Jj(-1)^j\zeta^{j-1}M_Z^j .
\]
Comparing coefficients of $n^{-s}$ for $n\le Z^J$ gives the identity
(cf.~\cite{HB82}).
\end{proof}

\begin{lemma}[Separation of variables]\label{lem:separation}
Let $k\ge1$ and let $F$ be smooth on $(0,\infty)^k$ and supported in
$[c,C]^k$.  Define
$F^\flat(\tau)=(2\pi)^{-k}\int_{(0,\infty)^k}F(s)\,s^{i\tau}\,ds/s$, where
$s^{i\tau}=\prod_j s_j^{i\tau_j}$ and $ds/s=\prod_j ds_j/s_j$.  Then for all
$t\in(0,\infty)^k$,
\[
 F(t)=\int_{\R^k}F^\flat(\tau)\,t^{-i\tau}\,d\tau ,
 \qquad
 |F^\flat(\tau)|\ll_{J,k,c,C}\|F\|_{C^{kJ}}\prod_{j=1}^k(1+|\tau_j|)^{-J}.
\]
In particular, if $\rho$ is smooth, supported in $[c/2,2C]$ and equal to
$1$ on $[c,C]$, then
$F(t)=\int F^\flat(\tau)\prod_j\rho(t_j)t_j^{-i\tau_j}\,d\tau$.
\end{lemma}

\begin{proof}
Substitute $s_j=\e^{u_j}$.  This reduces the claim to Fourier inversion for
the smooth compactly supported function $u\mapsto F(\e^{u})$, and the decay
estimate follows by repeated integration by parts.
\end{proof}

\section{A bilinear inequality of Korolev--Changa type}\label{sec:KC}

For a finite \emph{labelled} family $\mathcal Y$ of units modulo $q$ (residues
may repeat), let
\[
 I(q;\mathcal Y)=\#\bigl\{(y_1,y_2,y_3,y_4)\in\mathcal Y^4:\
 \inv{y_1}+\inv{y_2}\equiv\inv{y_3}+\inv{y_4}\ (\mathrm{mod}\ q)\bigr\}.
\]

\begin{lemma}\label{lem:KC}
Let $q\ge2$, let $c$ be a unit modulo $q$, let $\mathcal Y,\mathcal Z$ be
finite labelled families of units modulo $q$, and let $\gamma,\beta$ be
complex coefficients on them.  Then
\[
 \Bigl|\sum_{z\in\mathcal Z}\sum_{w\in\mathcal Y}
 \gamma_z\beta_w\,e_q(c\,\inv z\,\inv w)\Bigr|
 \le\|\gamma\|_2\|\beta\|_2\,
 \bigl\{q\,I(q;\mathcal Z)\,I(q;\mathcal Y)\bigr\}^{1/8}.
\]
\end{lemma}

This is the case $k=s=2$ of~\cite[Lemma~1]{KC}, whose proof is given
in~\cite{Kor18}.  There it is stated for interval supports, with the norms
$\|\gamma\|_{k/(k-1)}$ and $\|\beta\|_{s/(s-1)}$.  The proof uses the supports only to name the energies, and we
include it for completeness.

\begin{proof}
Call the left side $|\cB|$.  Write $S(z)=\sum_w\beta_we_q(c\inv z\inv w)$.
By Cauchy--Schwarz, $|\cB|^2\le\|\gamma\|_2^2\sum_z|S(z)|^2$.  Choose
$\rho_z$ with $|\rho_z|=1$ and $|S(z)|^2=\rho_z S(z)^2$.  Then
\[
 \sum_z|S(z)|^2=\sum_{w_1,w_2}\beta_{w_1}\beta_{w_2}
 G(\inv{w_1}+\inv{w_2}),
 \qquad G(\lambda)=\sum_z\rho_z e_q(c\lambda\inv z).
\]
Let $j(\lambda)$ be the number of pairs $(w_1,w_2)$ with
$\inv{w_1}+\inv{w_2}\equiv\lambda$.  Cauchy--Schwarz over $(w_1,w_2)$ gives
\[
 \sum_z|S(z)|^2\le\|\beta\|_2^2
 \Bigl(\sum_{\lambda\bmod q}j(\lambda)|G(\lambda)|^2\Bigr)^{1/2}.
\]
Cauchy--Schwarz over $\lambda$ bounds the last sum by
$(\sum_\lambda j(\lambda)^2)^{1/2}(\sum_\lambda|G(\lambda)|^4)^{1/2}$.  Here
$\sum_\lambda j(\lambda)^2=I(q;\mathcal Y)$.  Expanding $|G(\lambda)|^4$ and
using orthogonality over $\lambda$, with $c$ a unit, gives
$\sum_\lambda|G(\lambda)|^4\le q\,I(q;\mathcal Z)$.  Combining,
\[
 |\cB|^4\le\|\gamma\|_2^4\|\beta\|_2^4\,
 \bigl(q\,I(q;\mathcal Z)I(q;\mathcal Y)\bigr)^{1/2} .
\]
Taking fourth roots completes the proof.
\end{proof}

\section{Reciprocal energies averaged over the modulus}\label{sec:energy}

\begin{lemma}[Integers]\label{lem:energy-integers}
Let $N,Q\ge1$, let $\mathcal I\subset(N,CN]\cap\Z$, and for $q\ge1$ let
$I_{\mathcal I}(q)$ be the energy $I(q;\cdot)$ of the units of $\mathcal I$
modulo $q$.  Then for every $\eps>0$,
\[
 \sum_{Q<q\le2Q}I_{\mathcal I}(q)\ll_{C,\eps}(QN^2+N^4)(QN)^{\eps}.
\]
\end{lemma}

This is the case $k=2$ of~\cite[Lemma~2.3]{Irv14}, which generalizes
Fouvry and Shparlinski~\cite[Lemma~3]{FS11}.  We include the proof.

\begin{proof}
Clearing denominators, a counted quadruple satisfies $q\mid\Delta$, where
\[
 \Delta=(x_1+x_2)x_3x_4-(x_3+x_4)x_1x_2 .
\]
If $\Delta\ne0$, then $|\Delta|\ll N^3$, and such a quadruple contributes
for at most $\tau(|\Delta|)\ll N^{\eps}$ moduli; this gives $N^{4+\eps}$.
If $\Delta=0$, then $1/x_1+1/x_2=1/x_3+1/x_4$ in $\Q$.  Write the common
value as $a/b$ in lowest terms.  Each representation
$1/x+1/y=a/b$ satisfies $(ax-b)(ay-b)=b^2$, so there are at most
$2\tau(b^2)\ll N^{\eps}$ representations, since $b\ll N^2$.  There are
therefore $\ll N^{2+\eps}$ such quadruples, each counted for at most $Q+1$
moduli.
\end{proof}

We use the theorem of Bloom and Maynard in the following form.  Let
$\Q_{=u}$ be the set of reduced rationals in $[0,1]$ with denominator exactly
$u$, and let $\Q_{\le Q_0}=\bigcup_{u\le Q_0}\Q_{=u}$.  For a finite set
$\mathcal B\subset\Q$ let
$E_4(\mathcal B)=\#\{(b_1,\dots,b_4)\in\mathcal B^4:\ b_1+b_2=b_3+b_4\}$,
the equality being exact in $\Q$.

\begin{lemma}[{Bloom--Maynard~\cite[Theorem~2]{BM}}]\label{lem:BM}
There is an absolute constant $C_0>0$ with the following property.  Let
$Q_0\ge4$ and $n\ge1$, and let $\mathcal B\subset\Q_{\le Q_0}$ satisfy
$|\mathcal B\cap\Q_{=u}|\le n$ for every $u\le Q_0$.  Then
$E_4(\mathcal B)\le(\log Q_0)^{C_0}(Q_0n)^2$.
\end{lemma}

Theorem~2 of~\cite{BM} is stated for all energies $E_{2m}$, $m\ge2$, with
bound $(\log Q_0)^{C^m}(Q_0n)^m$; we use only $m=2$.

\begin{lemma}[Reduced fractions]\label{lem:energy-fractions}
Let $1\le M'\le U'$, and let $\cF$ be a set of reduced fractions $a/u$ with
$1\le|a|\le M'$ and $U'<u\le CU'$.  Then
$E_4(\cF)\ll_C(U'M')^2(\log 2U')^{C_0}$.  For $q\ge1$, let $I_\cF(q)$ count
quadruples in $\cF^4$ whose numerators and denominators are units modulo
$q$ and which satisfy
$a_1\inv{u_1}+a_2\inv{u_2}\equiv a_3\inv{u_3}+a_4\inv{u_4}\pmod q$.  Then
for every $\eps>0$ and $Q\ge1$,
\[
 \sum_{Q<q\le2Q}I_\cF(q)\ll_{C,\eps}
 \bigl(Q(U'M')^2+(U'M')^4\bigr)(QU')^{\eps}.
\]
\end{lemma}

\begin{proof}
Since $|a/u|<M'/U'\le1$, the set $\cF$ is covered by
$\cF_{-1}=\cF\cap[-1,0]$ and $\cF_0=\cF\cap[0,1]$.  Translation by an
integer preserves reduced denominators.  Each translate $\cF_k-k$ therefore
lies in $\Q_{\le Q_0}$, $Q_0=\max(4,CU')$, and has at most $M'$ elements
with any given exact denominator.  Lemma~\ref{lem:BM} gives
$E_4(\cF_k)\ll(U'M')^2(\log 2U')^{C_0}$, since energy is invariant under
translation.

For finite sets $\mathcal A_1,\dots,\mathcal A_4\subset\Q$, two
applications of Cauchy--Schwarz give the following, uniformly in
$t\in\Q$:
\[
 \#\{a_i\in\mathcal A_i:\ a_1+a_2=a_3+a_4+t\}
 \le\prod_i E_4(\mathcal A_i)^{1/4}.
\]
Summing over the $O(1)$ choices of pieces bounds $E_4(\cF)$.  For the second
claim, a counted quadruple satisfies $q\mid\Delta$ with
\[
 \Delta=a_1u_2u_3u_4+a_2u_1u_3u_4-a_3u_1u_2u_4-a_4u_1u_2u_3 .
\]
Quadruples with $\Delta=0$ are exact rational solutions; there are at most
$E_4(\cF)$ of them, each counted for at most $Q+1$ moduli.  Quadruples with
$\Delta\neq0$ number at most $|\cF|^4\ll_C(U'M')^4$, and each is counted for
at most $\tau(|\Delta|)\ll(U')^{\eps}$ moduli, since $0<|\Delta|\ll U'^4$.
\end{proof}

\section{An averaged bilinear estimate}\label{sec:bilinear}

\begin{proposition}\label{prop:bilinear}
Let $Q,W\ge1$ and $1\le M'\le U'$, and let $\cF$ be as in
Lemma~\ref{lem:energy-fractions}.  Suppose that for each
$q\in(Q,2Q]$ we are given:
\begin{itemize}
\item a unit $c_q$ modulo $q$;
\item a subset $\cF_q\subset\cF$ of fractions whose numerator and
denominator are units modulo $q$;
\item a subset $\cW_q\subset(W,CW]\cap\Z$ of units modulo $q$;
\item coefficients $\gamma_\rho$ on $\cF$ and $\beta_w$ on $(W,CW]$,
independent of $q$, with $\sum_\rho|\gamma_\rho|^2\le\Gamma^2$ and
$\sum_w|\beta_w|^2\le\mathrm B^2$.
\end{itemize}
For $\rho=a/u$ put
\[
 \cB_q=\sum_{\rho=a/u\in\cF_q}\ \sum_{w\in\cW_q}
 \gamma_\rho\beta_w\,e_q(c_q\,a\inv u\,\inv w).
\]
Then for every $\eps>0$,
\[
 \Bigl(\sum_{Q<q\le2Q}|\cB_q|^2\Bigr)^{1/2}
 \ll_{C,\eps}\Gamma\,\mathrm B\,Q^{3/8}(U'M')^{1/4}W^{1/4}
 \bigl\{(Q+(U'M')^2)(Q+W^2)\bigr\}^{1/8}(QU'W)^{\eps}.
\]
\end{proposition}

\begin{proof}
For each $q$, apply Lemma~\ref{lem:KC} with $\mathcal Z$ the family of
residues $z=u\inv a$, $\rho=a/u\in\cF_q$, so that $\inv z=a\inv u$, and
$\mathcal Y=\cW_q$.  The energy of $\mathcal Z$ is at most $I_\cF(q)$, and
that of $\cW_q$ at most $I_{\mathcal I}(q)$ with
$\mathcal I=(W,CW]\cap\Z$.  Squaring the bound of Lemma~\ref{lem:KC} and
summing over $q$, we apply H\"older's inequality with exponents $(4,4,2)$:
\[
 \sum_q|\cB_q|^2\le\Gamma^2\mathrm B^2(2Q)^{1/4}Q^{1/2}
 \Bigl(\sum_qI_\cF(q)\Bigr)^{1/4}\Bigl(\sum_qI_{\mathcal I}(q)\Bigr)^{1/4}.
\]
Lemmas~\ref{lem:energy-integers} and~\ref{lem:energy-fractions} complete the
proof.
\end{proof}

\section{A mean-value theorem for M\"obius--Kloosterman fractions}\label{sec:meanvalue}

Let $D\ge2$, $1\le E\le D$ and $H\ge1$.  Let $V$ be smooth and supported in
$[1,2]$.  Let $(\nu_h)$ be complex numbers supported on $1\le|h|\le H$ with
$|\nu_h|\le1$.  For $e\ge1$ put
\begin{equation}\label{eq:Pe}
 P_e=\sum_{(d,2e)=1}\mu(d)V(d/D)\sum_h\nu_h\,e_e(h\inv d),
 \qquad
 \cX=\Bigl(\sum_{E/2<e\le2E}|P_e|^2\Bigr)^{1/2}.
\end{equation}

\begin{theorem}\label{thm:meanvalue}
Assume $H\le D^{1/6}$ and $E\ge D^{2/3}H^2$.  Then for every $\eps>0$,
\[
 \cX\ll_\eps\|V\|_{C^4}\,D^{\eps}
 \Bigl(D^{11/12}E^{1/2}H^{3/4}+D^{1/3}EH+DH\Bigr).
\]
\end{theorem}

Note that the bound is uniform in the coefficients $\nu_h$.

\begin{proof}
We may assume that $D$ is large.  Throughout, $\eps$ denotes an arbitrarily
small positive constant, and polylogarithmic factors are absorbed into
$D^{\eps}$.

\emph{Step 1: gcd strata.}  For a frequency $h$ put $g=(h,e)$, $h=ga$ and
$e=gq$, so that $(a,q)=1$ and $e_e(h\inv d)=e_q(a\inv d)$.  Then
$P_e=\sum_{g\mid e}P^{(g)}_{e/g}$, where
\[
 P^{(g)}_q=\sum_{(d,2gq)=1}\mu(d)V(d/D)
 \sum_{\substack{1\le|a|\le H/g\\(a,q)=1}}\nu_{ga}\,e_q(a\inv d).
\]
By Minkowski's inequality in $\ell^2(e)$,
$\cX\le\sum_{g\le H}\cX_g$ with
$\cX_g=(\sum_{E/(2g)<q\le2E/g}|P^{(g)}_q|^2)^{1/2}$.  The range of $q$ is
the union of the two dyadic ranges $(Q,2Q]$ with $Q=E/(2g)$ and $Q=E/g$,
which we treat separately.  It therefore suffices to show that $\cX_g$ is
$g^{-1}$ times the bound of the theorem.

\emph{Step 2: Heath-Brown decomposition.}  Apply Lemma~\ref{lem:HB} with
$J=5$ and $Z=(4D)^{1/5}$ to $\mu(d)$ for $d\le2D$.  Decompose every atom
dyadically, using sharp dyadic intervals for the $m_i$ and a smooth dyadic
partition of unity for the $n_i$.  Write
$V(d/D)=\int V^\flat(\tau)(d/D)^{-i\tau}\,d\tau$ by
Lemma~\ref{lem:separation}, with $\int|V^\flat(\tau)|(1+|\tau|)^2\,d\tau\ll
\|V\|_{C^4}$.  The factor $(d/D)^{-i\tau}$ factors over the atoms.

Discard the $O((\log D)^9)$ boxes whose product range does not meet $[D,2D]$.
In every remaining box, the atoms have dyadic sizes whose product lies in
$[2^{-9}D,4D]$.  The coefficient attached to an atom has modulus at most
$1$, and $m_i\le Z\le D^{1/5+o(1)}$.

The following dichotomy holds for every box.  Either
\begin{enumerate}
\item[(B)] some set of atoms has product of sizes $P$ with
$D^{1/3}\le P<D^{2/3}$; or
\item[(L)] some $n_i$-atom has size $N>D^{2/3}$.
\end{enumerate}
Indeed, if neither some atom size lies in $[D^{1/3},D^{2/3})$ nor (L)
holds, then every atom size is below $D^{1/3}$.  The first partial product
exceeding $D^{1/3}$ is then below $D^{2/3}$.

\emph{Step 3: balanced boxes.}  In case (B), let $u$ be the product of the
atoms in the chosen set or in its complement, whichever has the smaller
size, and let $w$ be the product of the rest.  The size $U$ of $u$ satisfies
$2^{-9}D^{1/3}\le U\le2D^{1/2}$, and the size of $w$ is
$W\asymp D/U$.  Collapsing the atoms gives coefficients $\alpha_u,\beta_w$
with $|\alpha_u|\le\tau_9(u)$ and $|\beta_w|\le\tau_9(w)$.  The masks
$(uw,2gq)=1$ factor into $(u,2gq)=1$ and $(w,2gq)=1$, just as the mask
$(d,2gq)=1$ factors over the atoms.  For a fixed
$\tau$, the box contributes
\[
 \sum_{1\le|a|\le H/g}\ \sum_u\sum_w\nu_{ga}\alpha_u\beta_w\,
 e_q(a\inv u\,\inv w).
\]
Now stratify by $g'=(a,u)$: write $a=g'a_0$ and $u=g'u_0$, so that
$(a_0,u_0)=1$.  Since $g'$ is a unit modulo $q$,
$e_q(a\inv u\,\inv w)=e_q(a_0\inv{u_0}\,\inv w)$.  For fixed $g'$ the
fractions $\rho=a_0/u_0$ are distinct.  They satisfy $1\le|a_0|\le M'$ and
$U'<u_0\le CU'$ with
\[
 M'=\frac{H}{gg'},\qquad U'=\frac U{g'},
\]
and $M'\le U'$ because $H\le D^{1/6}<2^{-9}D^{1/3}\le U$.  The coefficient
$\gamma_\rho=\nu_{gg'a_0}\alpha_{g'u_0}$ satisfies
\[
 \sum_\rho|\gamma_\rho|^2\ll D^{\eps}U'M',
 \qquad
 \sum_w|\beta_w|^2\ll D^{\eps}W,
\]
by Lemma~\ref{lem:divisor}.  Proposition~\ref{prop:bilinear} with modulus
range $Q=E/g$ therefore bounds the $\ell^2(q)$ norm of this piece by
\[
 D^{\eps}(U'M'W)^{3/4}Q^{3/8}\{(Q+(U'M')^2)(Q+W^2)\}^{1/8}
 \le D^{\eps}g^{-5/4}g'^{-3/2}\,\Phi(U),
\]
where
\[
 \Phi(U)=(DH)^{3/4}E^{3/8}\{(E+(UH)^2)(E+W^2)\}^{1/8},
 \qquad W=D/U.
\]
The sum over $g'$ converges.  We claim that
$\Phi(U)\ll D^{11/12}E^{1/2}H^{3/4}$ for $2^{-9}D^{1/3}\le U\le2D^{1/2}$.
Since $U\le2D^{1/2}$, we have $W^2\asymp D^2/U^2\gg D\ge E$.  So only
$E\ll W^2$ can occur, and there are two cases.
\begin{enumerate}
\item If $(UH)^2\le E$, then
$\Phi\ll(DH)^{3/4}E^{1/2}W^{1/4}\ll D^{11/12}E^{1/2}H^{3/4}$, because
$W\ll D^{2/3}$.
\item If $E<(UH)^2$, then $\Phi\ll(DH)^{3/4}E^{3/8}(DH)^{1/4}
=DHE^{3/8}$.  This is $\le D^{11/12}E^{1/2}H^{3/4}$ exactly when
$E\ge D^{2/3}H^2$.
\end{enumerate}
(In Step~3 the hypothesis $H\le D^{1/6}$ is used only to ensure $M'\le U'$.)

\emph{Step 4: long boxes.}  In case (L), write $d=kn$, where $n$ is the long
smooth atom, of size $N>D^{2/3}$, and $k$ is the product of the other
atoms, of size $K\asymp D/N\ll D^{1/3}$.  The collapsed coefficient
$\alpha_k$ satisfies $\sum_k|\alpha_k|\ll D^{\eps}K$.  Fix $a$, $k$ and
$q$.  Detect $(n,2g)=1$ by $\sum_{\delta\mid(n,2g)}\mu(\delta)$; terms with
$(\delta,q)>1$ are empty.  Lemma~\ref{lem:incomplete}, applied to the
variable $v=n/\delta$ with $b=a\inv{k\delta}$ and $L\ll(1+|\tau|)^2$,
gives
\[
 \sum_{(n,2gq)=1}\rho_0(n/N)n^{-i\tau}e_q(a\inv k\,\inv n)
 \ll D^{\eps}(1+|\tau|)^2\bigl(q^{1/2}+Nq^{-1/2}\bigr).
\]
Summing absolutely over $k$ and over $|a|\le H/g$, and then taking the
$\ell^2$ norm over $q\in(E/g,2E/g]$, bounds the contribution of the box to
$\cX_g$ by
\[
 D^{\eps}(1+|\tau|)^2\frac Hg K\Bigl(\frac Eg+N\Bigr)
 \ll D^{\eps}(1+|\tau|)^2g^{-1}\bigl(D^{1/3}EH+DH\bigr).
\]

\emph{Step 5: conclusion.}  Integrate over $\tau$ against $|V^\flat(\tau)|$
and sum over the $O((\log D)^9)$ boxes and the $j\le5$ terms of
Lemma~\ref{lem:HB}.  Then sum over $g\le H$, with the factors
$g^{-5/4}$ and $g^{-1}$.  This proves the theorem.
\end{proof}

\section{Proof of Theorem~\ref{thm:main}}\label{sec:proof}

\subsection{Reduction to $x\le y$}
The substitution $(a,b,d,e)\mapsto(b,a,e,d)$ maps solutions of
$be-ad=\kappa$ to solutions of $be-ad=-\kappa$.  It interchanges the roles
of $(A,D)$ and $(B,E)$ and replaces $W$ by a permuted weight.  It therefore
suffices to treat $x\le y$, for both signs of $\kappa$.  Then $D\ge E$, and
we put
\[
 T=\frac EA=X^{1-x-y},\qquad
 \lambda=\frac{\log E}{\log D}=\frac{1-y}{1-x},\qquad
 \vartheta=\frac{\log T}{\log D}=\frac{1-x-y}{1-x}.
\]
Since $y\le1/2$, we have $\lambda+\vartheta\ge1$.  The condition
$8x+9y>8$ is equivalent to $\vartheta<1/9$.  On the compact set $K$ we
therefore have $\vartheta\le1/9-\eta$ for some $\eta=\eta(K)>0$.

\subsection{Poisson summation}
For odd coprime $d,e$, let $a_0=a_0(d,e)$ be the residue class modulo $2e$
determined by $a_0\equiv-\kappa\inv d\pmod e$ and $a_0\equiv1\pmod2$.  The
solutions of $be-ad=\kappa$ in odd integers with given $d,e$ are exactly
$a\equiv a_0\pmod{2e}$, $b=(ad+\kappa)/e$; for such $a$, $b$ is
automatically an odd integer.  Since $AD=BE=X$,
\[
 W\Bigl(\frac aA,\frac bB,\frac dD,\frac eE\Bigr)
 =F\Bigl(\frac aA;\frac dD,\frac eE\Bigr),
 \qquad
 F(\alpha;\delta,\epsilon)=W\Bigl(\alpha,\frac{\alpha\delta}{\epsilon}
 +\frac{\kappa}{\epsilon X},\delta,\epsilon\Bigr).
\]
Here $F$ is smooth, supported in $\alpha,\delta,\epsilon\in[1,2]$, and has
derivatives bounded independently of $X$.  Let
$\Phi(\xi;\delta,\epsilon)=\int F(\alpha;\delta,\epsilon)e(-\alpha\xi)\,d\alpha$.
Then for all $i,j,J\ge0$,
\begin{equation}\label{eq:Phi-decay}
 |\xi|^i\,\partial_\xi^i\partial_{\delta,\epsilon}^{j}
 \Phi(\xi;\delta,\epsilon)\ll_{i,j,J,W}(1+|\xi|)^{-J}.
\end{equation}
By CRT, $a_0/(2e)\equiv-\tfrac\kappa2\,\inv d/e+\tfrac12\pmod1$.  So
Lemma~\ref{lem:poisson} gives
\begin{equation}\label{eq:after-poisson}
 \cS_\kappa=\sum_{\substack{d,e\ \mathrm{odd}\\(d,e)=1}}\mu(d)\mu(e)
 \frac A{2e}\sum_{h\in\Z}(-1)^h\,
 \Phi\Bigl(\frac{hA}{2e};\frac dD,\frac eE\Bigr)\,
 e_e\Bigl(-\frac\kappa2\,h\inv d\Bigr).
\end{equation}

\subsection{The zero frequency}
The term $h=0$ is
\begin{equation}\label{eq:zero-mode}
 \cS_\kappa^0=\sum_{e\ \mathrm{odd}}\mu(e)\frac A{2e}
 \sum_{(d,2e)=1}\mu(d)\,\Phi\Bigl(0;\frac dD,\frac eE\Bigr).
\end{equation}
Here $x\le y\le1/2$, so $\sqrt X\le E\le D$.  For each odd $e\le2E$,
Lemma~\ref{lem:mobius} with $N=D$, $q=2e$ and
$G(\delta)=\Phi(0;\delta,e/E)$ bounds the inner sum by
$D\exp(-c\sqrt{\log D})$.
\begin{itemize}
\item The lemma applies because $q\le4E\le4D\le D^{100}$.
\item Its implied constant depends only on bounds for $G$ and $G'$ on a
fixed compact set.  By~\eqref{eq:Phi-decay}, these are uniform in $e$ and
$X$.
\end{itemize}
Summing over $e\in[E,2E]$ gives
$\cS_\kappa^0\ll AD\exp(-c\sqrt{\log D})\ll X\exp(-c'\sqrt{\log X})$,
since $D\ge\sqrt X$.

\subsection{Nonzero frequencies}
By~\eqref{eq:Phi-decay}, the terms with $|h|>X^{\eps}T$ contribute
$O(X^{-100})$.  Insert a smooth dyadic partition of unity in $|h|$ and treat
$h>0$ and $h<0$ separately.  On a block $|h|\asymp H\le X^{\eps}T$, the
function
\[
 (\eta,\delta,\epsilon)\mapsto\rho_0(\eta)\,
 \Phi\Bigl(\pm\eta\,\frac{HA}{2E\epsilon};\delta,\epsilon\Bigr),
 \qquad \eta=|h|/H,
\]
is smooth and supported in $[1/2,2]\times[1,2]^2$.  By~\eqref{eq:Phi-decay},
its $C^k$ norms are $\ll_{k}1$ uniformly in $H$.  Lemma~\ref{lem:separation}
writes it as an integral, against an $L^1$ density with rapid decay in
$\tau$, of tensors
\[
 \nu_h\,V_\tau(d/D)\,w_\tau(e/E),\qquad
 \nu_h=\mathbf 1_{\pm h>0}\,(-1)^h\rho(|h|/H)(|h|/H)^{-i\tau_1},
\]
with $V_\tau(\delta)=\rho(\delta)\delta^{-i\tau_2}$ and
$w_\tau(\epsilon)=\rho(\epsilon)\epsilon^{-i\tau_3}$.  Here $0\le\rho\le1$,
$|\nu_h|\le1$, $\nu$ is supported on $1\le|h|\le 4H$, and
$\|V_\tau\|_{C^4}\ll(1+|\tau|)^4$.  Replacing $h$ by $-\tfrac\kappa2h$,
which is $\pm h$, turns the phase in~\eqref{eq:after-poisson} into
$e_e(h\inv d)$ at the cost of relabelling $\nu$.  For one tensor, Cauchy--Schwarz over $e$ gives
\[
 \Bigl|\sum_{e}\mu(e)w_\tau(e/E)\frac A{2e}P_e\Bigr|
 \le\frac A{2E}(E+1)^{1/2}\cX,
\]
with $P_e$ and $\cX$ as in~\eqref{eq:Pe}, for $V=V_\tau$ and modes up to
$4H\le4X^{\eps}T$.

\subsection{Exponents}
On $K$ we have $\vartheta\le1/9-\eta$ and $\lambda\ge1-\vartheta$.  Hence
$4X^{\eps}T\le D^{1/6}$, and $E\ge D^{2/3}(4X^{\eps}T)^2$, because
$\lambda-\tfrac23-2\vartheta\ge\tfrac13-3\vartheta\ge3\eta$ (for $\eps$
small in terms of $\eta$).  Theorem~\ref{thm:meanvalue} therefore applies.
After integrating over $\tau$ and summing over the $O(\log X)$ blocks, it
gives
\[
 \cS_\kappa-\cS_\kappa^0\ll X^{1+3\eps}\Bigl(
 D^{-1/12}T^{3/4}+D^{-2/3}E^{1/2}T+E^{-1/2}T\Bigr).
\]
Indeed, $AE^{-1/2}\cdot D^{11/12}E^{1/2}T^{3/4}=X\cdot D^{-1/12}T^{3/4}$,
and similarly for the other two terms.  In powers of $X$, the three savings
are
\[
 \frac{8x+9y-8}{12}=\sigma,\qquad
 \frac{2x+9y-5}{6},\qquad
 \frac{2x+y-1}{2}.
\]
On $\Omega$ with $x\le y$ we have $x>7/16$, and then the second and third
savings exceed $\sigma$.  Indeed, their differences with $\sigma$ are
$(9y-4x-2)/12\ge(5x-2)/12>0$ and $(4x-3y+2)/12>0$.  This proves
\eqref{eq:main-power}.  Together with the bound for $\cS_\kappa^0$ it
proves \eqref{eq:main-log}, and Theorem~\ref{thm:main} follows.
Corollary~\ref{cor:diagonal} is the case $x=y$, where $8m+9M=17x$.
\qed

\section{Divisor pairs: proof of Theorem~\ref{thm:level}}\label{sec:level}

For $P\ge1$ and an integer $k\ge0$, we call a function $\phi\colon\R\to\C$
\emph{$(P,k)$-tempered} if it is smooth, supported in $[1/4,4]$, and
satisfies $\|\phi^{(j)}\|_\infty\le P^j$ for $0\le j\le k$.  For such $\phi$,
$k$ integrations by parts give $|\widehat\phi(\xi)|\le 8P^k(2\pi|\xi|)^{-k}$.

By splitting into $O(1)$ dyadic pieces, Theorem~\ref{thm:meanvalue}
applies, with a constant depending only on the support, to weights $V$
supported in any fixed compact subset of $(0,\infty)$ and to moduli $e$ in
any range $[E/c,cE]$.  Its constant depends polynomially on finitely many
derivatives of $V$.  The constant in Lemma~\ref{lem:separation} depends on
finitely many derivatives of $F$ in the same way.

\subsection{A smooth version}

\begin{proposition}\label{prop:smooth-level}
Let $0<\eta\le1/100$.  There is $k=k(\eta)$ such that the following holds.
Let $N_0\ge2$, let $1\le D,E\le N_0^{9/17-\eta}$, let $P\ge1$, and let
$\psi,\phi_1,\phi_2$ be $(P,k)$-tempered.  Then
\[
 \Delta:=\sum_{d,e}\mu(d)\mu(e)\phi_1\Bigl(\frac dD\Bigr)\phi_2\Bigl(\frac eE\Bigr)
 \Bigl[\sum_{\substack{d\mid n\\ e\mid n+2}}\psi\Bigl(\frac n{N_0}\Bigr)
 -N_0\widehat\psi(0)\varrho(d,e)\Bigr]
 \ll_\eta P^{k}N_0^{1-\eta}.
\]
\end{proposition}

\begin{proof}
Throughout, $\eps_0=\eta/48$, and $P^{O(1)}$ denotes a power of $P$ with an
exponent depending only on $\eta$.

\emph{Step 1: Poisson summation in $n$.}  Fix $d,e$ with
$\mu(d)\mu(e)\ne0$.  If $(d,e)\nmid2$, there is no $n$ and
$\varrho(d,e)=0$.  Otherwise the conditions $d\mid n$, $e\mid n+2$ define
one residue class $n_0$ modulo $q=[d,e]$.  Lemma~\ref{lem:poisson} gives
\[
 \sum_{n\equiv n_0\ (\mathrm{mod}\ q)}\psi\Bigl(\frac n{N_0}\Bigr)
 =\frac{N_0}q\sum_{h\in\Z}\widehat\psi\Bigl(\frac{hN_0}q\Bigr)
 e\Bigl(\frac{hn_0}q\Bigr).
\]
The term $h=0$ is $N_0\widehat\psi(0)\varrho(d,e)$, which cancels.

\emph{Step 2: the phases.}  Write $d=2^id'$ and $e=2^je'$ with $d',e'$ odd
and $i,j\in\{0,1\}$; then $(d',e')=1$.  Let $\inv{d'}$ be an inverse of $d'$
modulo $e'$.  Put $\ell=1$ and $c=2$ if $i=j=0$, and $\ell=2$ and $c=1$
otherwise.  Then $q=\ell d'e'$, and
\begin{equation}\label{eq:pair-phase}
 \frac{n_0}q\equiv-\frac{c\,\inv{d'}}{e'}\pmod 1 .
\end{equation}
For instance, if $i=j=0$, then $n_0=dt$ with $t\equiv-2\inv d\pmod e$, so
$n_0/(de)\equiv-2\inv d/e$.  If $i=0$ and $j=1$, then $n_0=dt$ with
$dt\equiv-2\pmod{2e'}$.  This forces $t=2t'$ with $t'\equiv-\inv d\pmod{e'}$,
so $n_0/(2de')\equiv-\inv d/e'$.  The cases $i=1$ are identical.

Since $\mu(2d')=-\mu(d')$ for odd $d'$, we get
$\Delta=\sum_{i,j}(-1)^{i+j}\Delta_{ij}$, where
\[
 \Delta_{ij}=\sum_{\substack{d',e'\ \mathrm{odd}\\(d',e')=1}}
 \mu(d')\mu(e')\phi_1\Bigl(\frac{d'}{D'}\Bigr)\phi_2\Bigl(\frac{e'}{E'}\Bigr)
 \frac{N_0}{\ell d'e'}\sum_{h\neq0}
 \widehat\psi\Bigl(\frac{hN_0}{\ell d'e'}\Bigr)
 e\Bigl(-\frac{c\,h\inv{d'}}{e'}\Bigr),
\]
with $D'=D/2^i$ and $E'=E/2^j$.

\emph{Step 3: orientation.}  By the reciprocity law
$\inv{d'}/e'+\inv{e'}/d'\equiv1/(d'e')\pmod 1$, where $\inv{e'}$ is an
inverse modulo $d'$, we have
$e(-ch\inv{d'}/e')=e(ch\inv{e'}/d')\,e(-ch/(d'e'))$.  On the relevant range
$|h|\le N_0^{\eps_0}T$, defined below, the second factor is a smooth
function of $(h/H,d'/D',e'/E')$ with bounded derivatives.  We may therefore
interchange the roles of $d'$ and $e'$ when $D'<E'$.  So assume $D'\ge E'$,
and put
\[
 T=\frac{\ell D'E'}{N_0},\qquad
 x=1-\frac{\log D'}{\log N_0},\qquad
 y=1-\frac{\log E'}{\log N_0},
\]
so that $x\le y$ and $x\ge8/17+\eta$.

\emph{Step 4: truncation.}  Since $|\widehat\psi(\xi)|\le8P^k(2\pi|\xi|)^{-k}$,
frequencies with $|h|>N_0^{\eps_0}T$ contribute $O(P^kN_0^{-10})$ once
$k\ge20/\eps_0$.  In particular, if $T\le N_0^{-\eps_0}$, then $\Delta_{ij}$
is negligible.  Assume from now on that $T>N_0^{-\eps_0}$, so that
$x+y\le1+2\eps_0$.

\emph{Step 5: separation and Cauchy--Schwarz.}  Insert a smooth dyadic
partition of unity in $|h|$, with blocks $|h|\asymp H\le N_0^{\eps_0}T$,
and treat each sign of $h$ separately.  Writing $\delta=d'/D'$ and
$\epsilon=e'/E'$, the weight is
\[
 \frac{N_0}{\ell d'e'}\widehat\psi\Bigl(\frac{hN_0}{\ell d'e'}\Bigr)
 \phi_1(\delta)\phi_2(\epsilon)
 =\frac1T\cdot\frac{1}{\delta\epsilon}\,
 \widehat\psi\Bigl(\pm\frac{|h|/H}{\delta\epsilon}\cdot\frac HT\Bigr)
 \phi_1(\delta)\phi_2(\epsilon).
\]
After multiplication by $\rho_0(|h|/H)$, it is $T^{-1}$ times a smooth
function of $(|h|/H,\delta,\epsilon)$ supported in a fixed box, with $C^m$
norms $\ll_m P^{O(m)}$.  Lemma~\ref{lem:separation} writes it as an integral
of tensors $\nu_h V_\tau(\delta)w_\tau(\epsilon)$, as in
Section~\ref{sec:proof}, against a density whose $L^1$ norm, weighted by
$(1+|\tau|)^4$, is $P^{O(1)}$.  Replacing $h$ by $-ch$ turns the phase into
$e_{e'}(h\inv{d'})$, with $\nu$ now supported on $|h|\le8H$ and still bounded
by $1$.

For one tensor, Cauchy--Schwarz over $e'\asymp E'$ bounds the contribution
by $T^{-1}(8E')^{1/2}\cX$, where $\cX$ is as in~\eqref{eq:Pe} with
$(D,E,H)$ replaced by $(D',E',8H)$.

\emph{Step 6: applying Theorem~\ref{thm:meanvalue}.}  Its hypotheses
$8H\le D'^{1/6}$ and $E'\ge D'^{2/3}(8H)^2$ read, in exponents,
\[
 5x+6y>5+O(\eps_0),\qquad 8x+3y>5+O(\eps_0).
\]
Both hold, since $y\ge x\ge8/17+\eta$ gives $5x+6y\ge11x>5.17$ and
$8x+3y\ge11x>5.17$.  Theorem~\ref{thm:meanvalue}, the identity
$T^{-1}=N_0/(\ell D'E')$ and $H\le N_0^{\eps_0}T$ give
\[
 \Delta_{ij}\ll P^{O(1)}N_0^{1+3\eps_0}\Bigl(D'^{-1/12}T^{3/4}
 +D'^{-2/3}E'^{1/2}T+E'^{-1/2}T\Bigr),
\]
as in Section~\ref{sec:proof}.  In powers of $N_0$, the three savings are
$\sigma_1=(8x+9y-8)/12$, $\sigma_2=(2x+9y-5)/6$ and $\sigma_3=(2x+y-1)/2$.
\begin{itemize}
\item Since $y\ge x$, we have $\sigma_1\ge(17x-8)/12\ge17\eta/12$.
\item $\sigma_2-\sigma_1=(9y-4x-2)/12\ge(5x-2)/12>0$.
\item Since $y\le1+2\eps_0-x$, we have
$\sigma_3-\sigma_1=(4x-3y+2)/12\ge(7x-1-6\eps_0)/12>0$.
\end{itemize}
Hence $\Delta_{ij}\ll P^{O(1)}N_0^{1-17\eta/12+3\eps_0}$.  Summing over the
$O(\log N_0)$ blocks, the two signs and the four parity cases gives the
claim, because $17\eta/12-4\eps_0\ge\eta$.
\end{proof}

\subsection{Proof of Theorem~\ref{thm:level}}
Let $\theta<9/17$, and fix $\eta\in(0,1/100]$ with $\theta\le9/17-3\eta$.
Let $k=k(\eta)$ be as in Proposition~\ref{prop:smooth-level}.  Put
$\delta'=\eta/(2k)$ and $\nu=N^{-\delta'}$.  Write $\Sigma(K,L)$ for the sum
in Theorem~\ref{thm:level} and
$C_N(d,e)=\#\{1\le n\le N:\ d\mid n,\ e\mid n+2\}$.

\emph{Trivial range.}  Since $|r_N(d,e)|\le1$, we have $|\Sigma(K,L)|\le KL$.
So we may assume $KL>N^{1-\delta'}$.  Then $K,L\ge N^{1-\delta'-\theta}\ge
N^{2/5}$, so that $\nu K,\nu L\ge N^{1/3}$.

\emph{Smoothing.}  Choose smooth functions $\chi_K,\chi_L$ and $\Psi$ with
the following properties.
\begin{itemize}
\item $\chi_K=1$ on $[1,(1-\nu)K]$ and $\chi_K=0$ on $[K,\infty)$, with
$\chi_K^{(j)}\ll_j(\nu K)^{-j}$; $\chi_L$ is defined in the same way.
\item $\Psi=1$ on $[2N^{1-\delta'},(1-\nu)N]$ and $\Psi=0$ outside
$[N^{1-\delta'},N]$, with $\Psi^{(j)}\ll_j(\nu N)^{-j}$.
\end{itemize}
Multiplying by a smooth dyadic partition of unity writes
\[
 \chi_K(d)=\sum_D\phi_D(d/D),\qquad
 \chi_L(e)=\sum_E\phi'_E(e/E),\qquad
 \Psi(n)=\sum_{N_0}\psi_{N_0}(n/N_0).
\]
These are sums of $O(\log N)$ functions, each $(P,k)$-tempered with
$P\ll\nu^{-1}$; here $D\le2K$, $E\le2L$ and $N^{1-\delta'}/2\le N_0\le N$.

Replacing $\mathbf 1_{d\le K}$ by $\chi_K(d)$ changes $\Sigma(K,L)$ by at most
$\sum_{d\in I}\sum_{e\le L}(C_N(d,e)+N\varrho(d,e))$, where
$I=((1-\nu)K,K]$.  By the divisor bound,
\[
 \sum_{d\in I}\sum_{e\le L}C_N(d,e)
 \le\sum_{n\le N}\tau(n+2)\,\#\{d\in I:\ d\mid n\}
 \ll N^{\eps}\sum_{d\in I}\frac Nd\ll N^{1+\eps}\nu.
\]
Since $\varrho(d,e)\le2(d,e)/(de)$, we have
$\sum_{e\le L}\varrho(d,e)\le2\tau(d)(1+\log L)/d$, so that
$N\sum_{d\in I}\sum_{e\le L}\varrho(d,e)\ll N^{1+\eps}\nu$.  The same holds
for $e$.

Next replace $\mathbf 1_{[1,N]}(n)$ by $\Psi(n)$ and $N$ by
$\int\Psi=\sum_{N_0}N_0\widehat{\psi_{N_0}}(0)$.  The two functions differ
only for $n\le2N^{1-\delta'}$ or $n\in[(1-\nu)N,N]$.  The first change
costs at most
$\sum_{n\text{ in these ranges}}\tau(n)\tau(n+2)\ll N^{1+\eps}(\nu+N^{-\delta'})$.
The second costs
$|N-\int\Psi|\sum_{d,e\le N}\varrho(d,e)\ll N^{1-\delta'}(\log N)^3$, using
$\sum_{d,e\le N}(d,e)/(de)\ll(\log N)^3$.

What remains is a sum of $O((\log N)^3)$ expressions $\Delta$ as in
Proposition~\ref{prop:smooth-level}.  Each has $D,E\le2N^\theta\le
N_0^{9/17-\eta}$, since $N_0\ge N^{1-\delta'}/2$ and $\delta'\le\eta$.  Each
is therefore $\ll\nu^{-k}N^{1-\eta}=N^{1-\eta/2}$.  Altogether
\[
 \Sigma(K,L)\ll N^{1-\eta/2}(\log N)^3+N^{1-\delta'+\eps},
\]
and Theorem~\ref{thm:level} follows with $\delta=\delta'/2$.
\qed

\section{Correlations of $\Lambda_R$: proof of Theorem~\ref{thm:LambdaR}}\label{sec:LambdaR}

Throughout this section $L(d)=\log(R/d)$, and
\[
 M(R)=\sum_{d,e\le R}\mu(d)\mu(e)L(d)L(e)\varrho(d,e).
\]

\begin{lemma}[Expansion]\label{lem:expansion}
For every integer $N\ge1$ and real $R$,
\[
 \sum_{n\le N}\Lambda_R(n)\Lambda_R(n+2)
 =N\,M(R)+\sum_{d,e\le R}\mu(d)\mu(e)L(d)L(e)\,r_N(d,e).
\]
\end{lemma}

\begin{proof}
Expand both factors and interchange summations.  The inner count is
$C_N(d,e)=N\varrho(d,e)+r_N(d,e)$.
\end{proof}

\begin{lemma}[Abel summation]\label{lem:abel}
Let $R\ge1$ and $K=\lfloor R\rfloor$, and let $a(d,e)$ be complex numbers for
$1\le d,e\le K$.  Suppose
$|\sum_{d\le k}\sum_{e\le l}a(d,e)|\le B$ for all $1\le k,l\le K$.  Then
\[
 \Bigl|\sum_{d,e\le K}a(d,e)L(d)L(e)\Bigr|\le B(\log R)^2.
\]
\end{lemma}

\begin{proof}
Let $w_k=\log((k+1)/k)$ for $1\le k<K$ and $w_K=\log(R/K)$.  Then
$w_k\ge0$, $L(d)=\sum_{k=d}^Kw_k$ for $1\le d\le K$, and
$\sum_{k=1}^Kw_k=\log R$.  Hence
\[
 \sum_{d,e\le K}a(d,e)L(d)L(e)
 =\sum_{k,l\le K}w_kw_l\sum_{d\le k}\sum_{e\le l}a(d,e),
\]
which is at most $B(\sum_kw_k)^2$ in absolute value.
\end{proof}

\subsection{The main term}
We use the prime number theorem for the M\"obius function with the
classical error term~\cite[Ch.~6]{MV}: for some $c_0>0$,
\begin{equation}\label{eq:PNTmu}
 M_\mu(x):=\sum_{n\le x}\mu(n)\ll x\exp\bigl(-c_0\sqrt{\log x}\bigr)
 \qquad(x\ge2).
\end{equation}
Below, $c_1,c_2,\dots$ denote positive absolute constants.

\begin{lemma}\label{lem:S1}
Put $m(t)=\sum_{n\le t}\mu(n)/n$ and
$S_1(z)=\sum_{n\le z}\mu(n)\log(z/n)/n$.  Then:
\begin{enumerate}
\item $m(t)\ll\exp(-c_1\sqrt{\log t})$ for $t\ge2$;
\item $S_1(z)=1+O(\exp(-c_2\sqrt{\log z}))$ for $z\ge1$.
\end{enumerate}
\end{lemma}

\begin{proof}
For real $\sigma>1$, partial summation gives
$\sum_n\mu(n)n^{-\sigma}=\sigma\int_1^\infty M_\mu(u)u^{-\sigma-1}\,du$.
This equals $1/\zeta(\sigma)$, which tends to $0$ as $\sigma\to1^+$.  By
\eqref{eq:PNTmu}, $|M_\mu(u)|u^{-2}$ is integrable on $[1,\infty)$, so
dominated convergence gives $\int_1^\infty M_\mu(u)u^{-2}\,du=0$.  Partial
summation then gives
\[
 m(t)=\frac{M_\mu(t)}t-\int_t^\infty\frac{M_\mu(u)}{u^2}\,du .
\]
Since
$\int_t^\infty\exp(-c_0\sqrt{\log u})\,du/u=\int_{\log t}^\infty
\e^{-c_0\sqrt v}\,dv\ll(1+\sqrt{\log t})\exp(-c_0\sqrt{\log t})$, this
proves~(1).

For~(2), $S_1(z)=\int_1^zm(t)\,dt/t$.  For $s>0$, absolute convergence
justifies interchanging sum and integral:
\[
 \int_1^\infty m(t)t^{-1-s}\,dt
 =\sum_n\frac{\mu(n)}n\int_n^\infty t^{-1-s}\,dt
 =\frac1{s\,\zeta(1+s)} .
\]
This tends to $1$ as $s\to0^+$.  By~(1) and dominated convergence,
$\int_1^\infty m(t)\,dt/t=1$.  Hence
$S_1(z)=1-\int_z^\infty m(t)\,dt/t$, which gives~(2) for $z\ge2$.  For
$1\le z<2$ we have $S_1(z)=\log z\in[0,1)$.
\end{proof}

\begin{lemma}\label{lem:Sg}
For $g\ge1$ and $y\ge2$, let
$S_g(y)=\sum_{d\le y,\ (d,g)=1}\mu(d)\log(y/d)/d$.  Then
\[
 S_g(y)=\frac g{\varphi(g)}+O\bigl(g^3\exp(-c_3\sqrt{\log y})\bigr).
\]
\end{lemma}

\begin{proof}
As in the proof of Lemma~\ref{lem:mobius}, we have
$\mu(n)\mathbf 1_{(n,g)=1}=\sum_{km=n,\ k\mid g^\infty}\mu(m)$, and hence
\[
 S_g(y)=\sum_{\substack{k\le y\\ k\mid g^\infty}}\frac{S_1(y/k)}k .
\]
Also $\sum_{k\mid g^\infty}1/k=\prod_{p\mid g}(1-1/p)^{-1}=g/\varphi(g)\le g$.
By Rankin's trick, for $Y\ge1$,
\[
 \sum_{\substack{k\mid g^\infty\\ k>Y}}\frac1k\le
 Y^{-1/2}\prod_{p\mid g}(1-p^{-1/2})^{-1}\le Y^{-1/2}4^{\omega(g)}\le
 Y^{-1/2}g^2 .
\]
Write
$S_g(y)-g/\varphi(g)=\sum_{k\le y,\,k\mid g^\infty}(S_1(y/k)-1)/k
-\sum_{k\mid g^\infty,\,k>y}1/k$.  Use Lemma~\ref{lem:S1}(2) for
$k\le\sqrt y$, and the bound $|S_1-1|\le1+\log2$ together with Rankin's
trick for $k>\sqrt y$.  This gives
$|S_g(y)-g/\varphi(g)|\ll g\exp(-c_2\sqrt{\tfrac12\log y})+g^2y^{-1/4}$,
which is acceptable.
\end{proof}

\begin{lemma}[gcd expansion]\label{lem:gcd}
Let $h$ be the multiplicative function on squarefree integers with
$h(2)=1$ and $h(p)=-1$ for odd primes $p$.  For squarefree $d,e$,
$\varrho(d,e)=(de)^{-1}\sum_{g\mid(d,e)}h(g)$.  Consequently
\[
 M(R)=\sum_{g\le R}\frac{\mu(g)^2h(g)}{g^2}\,S_g(R/g)^2 .
\]
\end{lemma}

\begin{proof}
The sum $\sum_{g\mid(d,e)}h(g)=\prod_{p\mid(d,e)}(1+h(p))$ equals
$(d,e)$ if $(d,e)\mid2$ and $0$ otherwise.  This is $de\,\varrho(d,e)$.
Inserting this into $M(R)$ and writing $d=gd'$, $e=ge'$, we use
$\mu(gd')=\mu(g)\mu(d')\mathbf 1_{(g,d')=1}$ and
$L(gd')=\log((R/g)/d')$.  This gives the formula.
\end{proof}

\begin{proposition}\label{prop:mainterm}
For $R\ge2$, $M(R)=\mathfrak S+O(\exp(-c_4\sqrt{\log R}))$.
\end{proposition}

\begin{proof}
The function $f(g)=\mu(g)^2h(g)/\varphi(g)^2$ is multiplicative and
supported on squarefree $g$.  It satisfies $f(2)=1$ and
$f(p)=-1/(p-1)^2$ for odd $p$, and $\sum_g|f(g)|<\infty$.  Hence
$\sum_gf(g)=\prod_p(1+f(p))=\mathfrak S$.  Let $G=\exp(\beta\sqrt{\log R})$
with $\beta=c_3/10$, and assume that $R$ is large; for bounded $R$ the claim
is trivial.  By Lemma~\ref{lem:gcd},
\[
 M(R)-\mathfrak S=
 \sum_{g\le G}\frac{\mu^2h}{g^2}\Bigl(S_g(R/g)^2-\frac{g^2}{\varphi(g)^2}\Bigr)
 +\sum_{G<g\le R}\frac{\mu^2h}{g^2}S_g(R/g)^2
 -\sum_{g>G}f(g).
\]
We bound the three sums in turn.
\begin{itemize}
\item For $g\le G$ we have $R/g\ge2$, and Lemma~\ref{lem:Sg} gives
$|S_g^2-g^2/\varphi(g)^2|\ll g^6\exp(-c_3\sqrt{\log(R/g)})$.  Since
$\sqrt{u-\beta\sqrt u}\ge\sqrt u-\beta$, the first sum is
$\ll G^5\exp(-c_3\sqrt{\log R})\le\exp(-\tfrac{c_3}2\sqrt{\log R})$.
\item The trivial bound $|S_g(y)|\le(1+\log y)^2$ shows that the second sum
is $\ll(\log R)^4/G$.
\item By Rankin's trick, the third sum is at most
$G^{-1/2}\sum_g\mu(g)^2g^{1/2}/\varphi(g)^2\ll G^{-1/2}$.
\end{itemize}
\end{proof}

\subsection{Proof of Theorem~\ref{thm:LambdaR}}
Let $\theta<9/17$, $N\ge2$ and $2\le R\le N^\theta$.  By
Lemma~\ref{lem:expansion},
\[
 \sum_{n\le N}\Lambda_R(n)\Lambda_R(n+2)-\mathfrak SN
 =N\bigl(M(R)-\mathfrak S\bigr)
 +\sum_{d,e\le R}\mu(d)\mu(e)L(d)L(e)\,r_N(d,e).
\]
Proposition~\ref{prop:mainterm} bounds the first term by
$N\exp(-c_4\sqrt{\log R})$.  For the second, apply Lemma~\ref{lem:abel} with
$a(d,e)=\mu(d)\mu(e)r_N(d,e)$.  The box sums are bounded by
Theorem~\ref{thm:level}, since $k,l\le\lfloor R\rfloor\le N^\theta$, so the
second term is at most $C_\theta N^{1-\delta}(\log R)^2$.  Finally, $\log
R\le\theta\log N$ gives
$N^{-\delta}(\log R)^2\le R^{-\delta/\theta}(\log R)^2\ll_\theta
\exp(-c_4\sqrt{\log R})$.
\qed

\section{The literature bar and the barrier at level $4/7$}\label{sec:barrier}

This section records what we know about Theorem~\ref{thm:level} beyond the
exponent $9/17$.  It is a synthesis, together with a report on an attempt to
pass level $4/7$ that did not succeed.  Its status differs from that of the
rest of the paper.
\begin{itemize}
\item A statement called a \emph{Claim} is an exponent computation: a
published estimate is inserted into a reduction and the conditions are
compared.  The Claims were produced by AI systems, have not been checked by
a human, and are not formalized.  The hypotheses of the cited estimates have
not been verified in detail.
\item A statement called a \emph{Lemma} is elementary and is proved here.
The Lemmas, and all exponent arithmetic of this section, are machine-checked
(\texttt{Barrier.lean}; see Section~\ref{sec:lean}).
\item The numerical experiments are finite computations.  The programs are
in the repository (\texttt{notes/lab/}).  They are evidence about what is
true, not about what can be proved.
\end{itemize}
Working notes with more detail are in \texttt{notes/} in the repository.

\emph{Units.}  Sizes are powers of $N$.  For a level $\theta$ we write
\[
 q\asymp N^{\theta},\qquad L=N^{\theta/2},\qquad H=N^{2\theta-1},\qquad
 K=N^{1-3\theta/2}.
\]
Here $q$ is a modulus, $H=q^2/N$ is the length of the dual sum after Poisson
summation, and $K$ is the mean of the count in~\eqref{eq:nu} below.  At
$\theta=4/7$ one has $H=K=L^{1/2}$.  All variables carry smooth dyadic
weights, which we suppress.

\subsection{From a saving to a level}\label{sec:dictionary}
After Poisson summation in $n$ (Section~\ref{sec:level}),
Theorem~\ref{thm:level} at level $\theta$ is a bound
\begin{equation}\label{eq:barrier-T}
 \sum_{0<|h|\le H}\ \sum_{d\asymp q}\ \sum_{e\asymp q}\alpha_d\lambda_e\,
 e\Bigl(\frac{2h\inv d}{e}\Bigr)\ll q^{2}N^{-\delta},
\end{equation}
with coefficients $\alpha,\lambda$ of M\"obius type.  The trivial bound is
$Hq^2$, so a saving of $HN^{\delta}$ is needed.

\begin{lemma}[Dictionary]\label{lem:dictionary}
Let $\kappa<2$.  A saving of $q^{\kappa}$ over the trivial bound
in~\eqref{eq:barrier-T} exceeds $H$ exactly when $\theta<1/(2-\kappa)$.
\end{lemma}
\begin{proof}
$H<q^{\kappa}$ means $2\theta-1<\kappa\theta$.
\end{proof}

\medskip
\begin{center}
\begin{tabular}{lll}
$\kappa$ & level $\theta$ & source \\ \hline
$1/48$ & $48/95$ & Duke--Friedlander--Iwaniec~\cite{DFI}, as in~\cite{CL}\\
$1/17$ & $17/33$ & Bettin--Chandee~\cite{BC}; Appendix~\ref{app:BC}\\
$1/9$ & $9/17$ & Theorem~\ref{thm:level};~\cite{BCHB} for mollifiers\\
$1/8$ & $8/15$ & Irving~\cite{Irv14}, transposed; Remark~\ref{rem:comparison}\\
$1/4$ & $4/7$ & Conrey~\cite{Con89} with~\cite{DI82}, transposed; Claim~\ref{claim:conrey}\\
$1/3$ & $3/5$ & Maynard~\cite{May2}, given a flexible factor; Claim~\ref{claim:window}\\
$2/5$ & $5/8$ & Pascadi~\cite{Pas25}, given a flexible factor; Claim~\ref{claim:window}
\end{tabular}
\end{center}
\medskip

\begin{claim}[The bar]\label{claim:conrey}
For $K,L\le R$ the sum in Theorem~\ref{thm:level} is
$\ll N^{\eps}(R^{7/4}+N^{1/2}R^{7/8})$.  In particular
Theorems~\ref{thm:LambdaR} and~\ref{thm:level} hold for every
$\theta<4/7$.
\end{claim}

\emph{Derivation.}  Conrey's treatment of the mollified second moment, in
the form given by Pratt and Robles~\cite[\S3.2.6]{PR18}, bounds a trilinear
form $\sum_a\nu(a)\sum_{u,v}\mu^\flat(u)r(v)e(-a\inv u/v)$ with $\nu$ and $r$
arbitrary, $\mu^\flat$ of M\"obius type, and $a$ of length $UV/T$.  This is
\eqref{eq:barrier-T}, with $H=DE/N$ in place of $UV/T$.  The Heath-Brown
identity is applied to $\mu^\flat$; sums with a long smooth variable are
treated by Weil's bound, and the remaining bilinear sums by the estimate of
Deshouillers and Iwaniec~\cite{DI82}, which takes arbitrary coefficients
inside and absolute values over the modulus.  Both terms of the resulting
bound are $N^{1-\delta}$ exactly for $\theta<4/7$.

\subsection{Using the structure of both variables}\label{sec:twosided}
Claim~\ref{claim:conrey} uses the arithmetic of one variable only.  It is
uniform in the shift, which is why it also serves the mollifier problem.
For a fixed shift more is available.

\emph{A fixed shift is a fixed residue class.}  The sum
in~\eqref{eq:barrier-T} is the Poisson dual of
\[
 \sum_{e}\lambda_e\Bigl(\#\{n=ad\le N:\ n\equiv-2\ (\mathrm{mod}\ e)\}
 -\text{expected}\Bigr),
\]
with $a$ smooth of length $N^{1-\theta}$.  This is the setting of Bombieri,
Friedlander and Iwaniec~\cite{BFI}: a fixed residue class and factorable
moduli.  Estimates of that kind are not uniform in the residue class and so
are not available for mollifiers.

\begin{claim}[One flexible side suffices]\label{claim:window}
Suppose that $\lambda_e$ is a combination of products $e=rs$ in which $s$
has size $N^{\sigma}$ with
\[
 2\theta-1<\sigma<\min\Bigl(\frac{2(1-\theta)}3,\ 2-3\theta\Bigr).
\]
Then~\eqref{eq:barrier-T} holds for arbitrary bounded coefficients
$\alpha_d$.  The interval is nonempty for $\theta<3/5$.  Without its last
term it is nonempty for $\theta<5/8$.
\end{claim}

\emph{Derivation.}  Apply Cauchy--Schwarz over $(d,r)$ and then
\cite[Theorem~12]{DI82} with modulus variable $r$, level $ss'$ and
frequencies $hs'-h's$.  The lower bound on $\sigma$ comes from the diagonal.
The bound $2(1-\theta)/3$ comes from the regular spectrum.  The bound
$2-3\theta$ comes from the exceptional spectrum; Pascadi's
refinement~\cite{Pas25} removes it for sequences of this structure.  The
resulting exponents agree with the levels $3/5$ of Maynard~\cite{May2} and
$5/8$ of Pascadi for triply well-factorable weights.  By the symmetry of
$be-ad=2$ the roles of $d$ and $e$ can be exchanged.

\medskip
Now decompose both $\mu(d)$ and $\mu(e)$ by the Heath-Brown identity.  A
piece is covered by Claim~\ref{claim:window} as soon as some sub-product of
the factors of $d$, or of $e$, has size in the window.  Pieces with a smooth
factor longer than $N^{5\theta/2-1}$ are covered by Weil's bound.

\begin{lemma}\label{lem:balanced-window}
For $0<\theta<1$, the size $\theta/2$ lies in
$[2\theta-1,\min(2(1-\theta)/3,2-3\theta)]$ if and only if
$\theta\le4/7$.  The same holds without the term $2-3\theta$.
\end{lemma}
\begin{proof}
$\theta/2\le2-3\theta$ and $\theta/2\le2(1-\theta)/3$ are each equivalent
to $\theta\le4/7$, and $2\theta-1\le\theta/2$ holds for $\theta\le2/3$.
\end{proof}

So for $\theta>4/7$ the balanced factorization $e=\ell_1\ell_2$,
$\ell_i\asymp N^{\theta/2}$, with both factors smooth, is not covered, and
removing the exceptional spectrum does not help it.  The Heath-Brown
identity contains this piece with a nonzero coefficient.  The uncovered
pieces are, up to factors of size $N^{O(\theta-4/7)}$:
\begin{itemize}
\item $e=\ell_1\ell_2$ with both factors smooth and of size near
$N^{\theta/2}$;
\item $e=\ell_1\ell_2$ with both factors smooth and of sizes near
$N^{1-\theta}$ and $N^{2\theta-1}$.
\end{itemize}
When $d$ and $e$ are both of the first kind, all six variables are smooth.

\medskip
\noindent\textbf{The missing estimate (T).}  \emph{For some $\theta>4/7$
and some $\delta>0$, the number of solutions of
\begin{equation}\label{eq:T}
 b\,\ell_1\ell_2-a\,n_1n_2=2,\qquad a,b\asymp N^{1-\theta},\quad
 \ell_1,\ell_2,n_1,n_2\asymp L,
\end{equation}
counted with smooth weights, equals its expected main term up to a factor
$1+O(N^{-\delta})$.}

\medskip
This is a restricted case of the additive divisor problem for
$d_3\times d_3$.  At $\theta=4/7$ the box is $(3/7,2/7,2/7)$ on each side.
An asymptotic formula for $\sum_{n\le x}d_3(n)d_3(n+h)$ is not known; the
cube $(1/3,1/3,1/3)$ corresponds to $\theta=2/3$.  Equivalently, (T) asks
for a level of distribution beyond $4/7$ for a restricted $d_3$, in a fixed
residue class, on average over moduli $\ell_1\ell_2$ with two smooth
factors.  For comparison:
\begin{itemize}
\item for $d_3$ and individual moduli the known levels are $1/2+1/82$
(Heath-Brown~\cite{HB86}), $1/2+1/46$ for prime moduli
(Fouvry--Kowalski--Michel~\cite{FKM3}) and $1/2+1/30$
(Sharma~\cite{Sha});
\item on average over moduli with a fixed residue class, \cite{FKM3} reach
$9/17$;
\item for moduli with a suitable factorization, the Type~III argument of
Polymath8a~\cite{Polymath8a} reaches $4/7$ (Section~\ref{sec:why47}).
\end{itemize}

\begin{claim}[Conditional and partial results]\label{claim:conditional}
\leavevmode
\begin{itemize}
\item If (T) and its analogues for the other uncovered pieces hold at
level $\theta$, then Theorems~\ref{thm:LambdaR} and~\ref{thm:level} hold at
that level, for $\theta<3/5$.
\item Unconditionally, the analogue of Theorem~\ref{thm:level} for weights
with no uncovered piece on at least one side holds for $\theta<3/5$.
Examples are well-factorable sieve weights, or truncated divisor sums over
$N^{\eta}$-smooth divisors.
\item For the true $\Lambda_R$, bounding the uncovered pieces trivially
extends Claim~\ref{claim:conrey} only to
$R\le N^{4/7}\exp\bigl(o((\log N)^{2/3})\bigr)$.
\end{itemize}
\end{claim}

The pieces that mix the two kinds above are a $d_3\times d_2$ problem with a
linear twist of size $N^{1/7}$, and a quadratic divisor problem with
coefficients of size $N^{1/7}$.  They may be within reach of existing
spectral results; we have not examined them.

\subsection{Equivalent forms of the missing estimate}\label{sec:forms}
We list the forms of (T) that we found.  Each is obtained from the previous
ones by Poisson summation, reciprocity, Cauchy--Schwarz or a spectral
expansion.  In every form the known method gives exactly the bound that is
needed with $\delta=0$.

\medskip
\noindent\textbf{(F1) Dual form.}
\[
 \sum_{0<|h|\le H}\ \sum_{n_1,n_2,\ell_1,\ell_2\asymp L}
 e\Bigl(\frac{2h\,\inv{n_1n_2}}{\ell_1\ell_2}\Bigr)\ll L^{4}N^{-\delta}.
\]
The trivial bound is $HL^4$.  By reciprocity
(Lemma~\ref{lem:roots}\,(iii)) the two sides $n_1n_2$ and $\ell_1\ell_2$
can be exchanged.

\medskip
\noindent\textbf{(F2) Hyper-Kloosterman form.}  Poisson summation in $n_1$
and $n_2$ gives
\[
 \sum_{h}\ \sum_{u,v\asymp L}\ \sum_{\ell_1,\ell_2\asymp L}
 \mathrm{Kl}_3(2huv;\ell_1\ell_2),
\]
and a saving of $H^{1+\delta}$ over Deligne's bound is needed.  Every
balanced variable has length equal to the square root of the modulus, so it
is self-dual under Poisson summation.

\medskip
\noindent\textbf{(F3) Quadratic divisor problem.}  For coefficients
$\ell_1,n_1\asymp L$ put
\[
 \mathcal D(\ell_1,n_1)=\sum_{\ell_1X-n_1Y=2}d^\flat(X)d^\flat(Y),
\]
where $X=b\ell_2$, $Y=an_2$ and $d^\flat$ counts factorizations in the
prescribed ranges.  At $\theta=4/7$ the main term has size $L^{3/2}$.  The
spectral method gives each $\mathcal D(\ell_1,n_1)$ with an error of size
$L^{7/4}$.  Needed: an error $L^{3/2-\delta}$ on average over
$(\ell_1,n_1)$.

\medskip
\noindent\textbf{(F4) Variance of the divisor function in thin classes.}
For $d\asymp L^2$ and $r\asymp L$ let
\begin{equation}\label{eq:nu}
 \nu(d,r)=\#\{(a,b,s):\ brs-ad=2,\ s\asymp L,\ a\asymp N^{1-\theta}\}
 =\sum_{k\le K}d_L(m_0+kd),
\end{equation}
where $d_L$ counts divisors in $[L,2L)$ and $m_0\equiv2\inv r$ modulo $d$.
The mean of $\nu$ is of size $K$.  By Cauchy--Schwarz over $(d,r)$, (T)
follows from
\[
 \operatorname{Var}_{(d,r)}\nu\ll K^{2-\delta}.
\]
The method of Deshouillers and Iwaniec gives $K^{2+\eps}$ at $\theta=4/7$.
In words: the divisor function in progressions modulo $d$, at length
$d^{5/4}$, has small variance over the thin set of classes
$\{2\inv r:\ r\le\sqrt d\}$.  Over all classes this is elementary, because
two solutions with the same $(r,d)$ satisfy
\[
 (bs-b's')\,r=(a-a')\,d ,
\]
so the pairs are parametrized by $bs-b's'=kd$; the difficulty is the
condition that $2\,\inv{bs}$ has a representative of size $L$ modulo $d$.
Bounding the variance over the thin set by the variance over all classes
loses a factor $L^{1/2}$.

\medskip
\noindent\textbf{(F5) A mean square at one level.}  Expanding the variance
in (F4), for a pair $s,s'\asymp L$ with $q=ss'$ one meets
\begin{equation}\label{eq:G}
 G_q(n)=\sum_{r\asymp L}\ \sum_{d\asymp L^2}
 e\Bigl(\frac{2n\,\inv d}{rq}\Bigr),
\end{equation}
evaluated at the frequencies $n=hs'-h's$ with $0<|h|,|h'|\le H$.  These
frequencies form a Bohr set $\mathcal N(s,s')$ of density $L^{-1/2}$ in
$[-2HL,2HL]$.  Consider the full mean square
\[
 S(q)=\sum_{0<|n|\le 2HL}|G_q(n)|^2 .
\]
\begin{center}
\begin{tabular}{ll}
bound for $S(q)$ at $\theta=4/7$ & size \\ \hline
Weil's bound, term by term & $L^{13/2}$\\
Kuznetsov's formula and the dual large sieve~\cite{DI82} & $L^{5+\eps}$\\
needed & $L^{5-\delta}$\\
observed (Section~\ref{sec:numerics}) & $L^{9/2}$
\end{tabular}
\end{center}

\noindent\textbf{(LS).}  \emph{$S(q)\ll L^{5-2\delta}$ for all but a
negligible set of $q=ss'$.}

\begin{claim}\label{claim:LS}
(LS) implies (T) at $\theta=4/7$, and the corresponding statement for
$\theta$ slightly larger implies (T) there.
\end{claim}

\emph{Derivation.}  Let $W$ be the sum in (F1), written as a sum over
$r,s\asymp L$ and $d=n_1n_2$.  Cauchy--Schwarz over $(r,d)$ removes the
divisor weight on $d$ and gives
\[
 |W|^2\le L^{3+\eps}\sum_{s,s'}\ \sum_{h,h'}G_{ss'}(hs'-h's).
\]
The terms with $hs'=h's$ contribute $L^{15/2}$ to the right-hand side,
which is admissible.  For the others, the terms of $S(q)$ are non-negative
and $\mathcal N(s,s')$ has $\asymp L$ elements, so
\[
 \Bigl|\sum_{n\in\mathcal N(s,s')}G_q(n)\Bigr|
 \le L^{1/2}\Bigl(\sum_{n\in\mathcal N(s,s')}|G_q(n)|^2\Bigr)^{1/2}
 \le L^{1/2}S(q)^{1/2}.
\]
With (LS) this is $L^{3-\delta}$, the sum over the $L^2$ pairs is
$L^{5-\delta}$, and $W\ll L^{4-\delta/2}$.  The size $L^{5+\eps}$ in the
table is our computation from the large sieve; the contribution of the
exceptional spectrum has not been rechecked in this form.

\medskip
The point of (F5) is that the structure of the Bohr set is not needed.  The
loss in the method of Deshouillers and Iwaniec is not the passage from the
sparse set to all frequencies; it is already present in the bound for
$S(q)$.  By the same positivity the modulus need not be a product: a bound
$\sum_qS(q)\ll L^{7-\delta}$ over all integers $q\asymp L^2$ implies (LS)
for most $q=ss'$.

The mean square $S(q)$ is a count of close pairs of the fractions
in~\eqref{eq:G}.  If $xk-cqz=1$ and $x'k'-c'qz'=1$, and $J=xc'-x'c$,
$\tilde\jmath=zk'-z'k$, then
\[
 Jkk'-\tilde\jmath\,cc'q=k'c'-kc,\qquad
 uu'+1=J\tilde\jmath\,q\quad(u=c'qz-x'k,\ u'=xk'-cqz').
\]
Here $J/(cc'q)$ is the difference of the fractions $\inv k/(cq)$ and
$\inv{k'}/(c'q)$, and $S(q)$ counts the solutions with $|J|\le L^{5/2}$.
These are the entries of an element of $\Gamma_0(q)$, so this is the usual
unfolding and leads back to Kuznetsov's formula.

\medskip
\noindent\textbf{(F6) Bessel's inequality for one vector.}  Let $g$ run over
an orthonormal basis of automorphic forms of level $q$ in a bounded spectral
window, with Fourier coefficients $\rho_g(n)$ and spectral parameters $t_g$.
Put $v_n=(\rho_g(n))_g$ and $a=\sum_{m\asymp\sqrt q}v_m$.  By Kuznetsov's
formula,
\[
 G_q(n)\ \propto\ \langle D_Xa,\,v_{2n}\rangle,\qquad
 (D_Xc)_g=X^{2it_g}c_g,\qquad X\approx q^{-7/8}.
\]
For $n\le q^{3/4}$ the vectors $v_n$ are almost orthonormal.  Let $V$ be
their span, of dimension $q^{3/4}$ in a space of dimension $\asymp q$.  Then
$S(q)\propto\|\Pi_VD_Xa\|^2$, and:
\begin{itemize}
\item Bessel's inequality $\|\Pi_VD_Xa\|^2\le\|a\|^2$ is the bound
$L^{5+\eps}$, that is, level $4/7$;
\item (LS) is the statement $\|\Pi_VD_Xa\|^2\le q^{-\delta}\|a\|^2$;
\item a vector in general position would give $q^{-1/4}\|a\|^2$, which is
what the experiments show.
\end{itemize}
A heuristic reading, which we have not made precise: $D_X$ is the wave
propagator on $\Gamma_0(q)\backslash\mathbb H$ at time $2\log(1/X)\approx
\tfrac74\log q$, the vector $a$ is a bump at the Fricke point $i/\sqrt q$,
and $V$ consists of the first $q^{3/4}$ Fourier modes on a closed horocycle.
At that time the wavefront covers the surface about $q^{3/4}$ times.
Weil's bound controls each sheet; nothing known controls the interference
between sheets.

\medskip
\noindent\textbf{(F7) A second moment with an Atkin--Lehner sign.}  Without
Cauchy--Schwarz, apply Kuznetsov's formula in $(\ell_1,n_1)$ at level
$q=\ell_2n_2$ for the pair of cusps $(\infty,1/\ell_2)$.  The dual sum of
(F1) becomes
\[
 \sum_{\ell_2,n_2\asymp L}\ \sum_{f\in\mathcal B(\ell_2n_2)}
 \omega_f\,h_X(t_f)\,\eta_f(n_2)\,L(\tfrac12,f)
 \sum_{h\asymp q^{1/4}}\lambda_f(h),
\]
where $\eta_f(n_2)$ is the Atkin--Lehner eigenvalue and $h_X$ oscillates
like $X^{2it}$.  The sum over $h$ is the dual half of an unbalanced
approximate functional equation.  So this is a second moment of
$L(\frac12,f)$ over the level $\ell_2n_2$, twisted by a partial
Atkin--Lehner sign.  It is the spectral form of (F3).
\begin{itemize}
\item The trivial bound, a second moment of size $\log q$, is the error
$L^{7/4}$ in (F3).
\item Needed: a saving of $q^{1/8+\delta}$ on average over
$(\ell_2,n_2)$.
\item If $n_2$ were absent from the level, the same argument would succeed
with room $L^{1/4}$.  It tolerates $n_2\le L^{1/2}$; the problem has
$n_2\asymp L$.
\end{itemize}

\medskip
\noindent\textbf{(F8) Roots of a quadratic congruence.}

\begin{lemma}\label{lem:roots}
Let $q,c,d$ be integers with $(d,cq)=1$ and put $D=qc^2+d^2$.  Let $x$ and
$e$ be integers with $dx\equiv1$ modulo $cq$ and $de\equiv1$ modulo $D$,
and put $\nu=ce$.  Then
\begin{itemize}
\item[(i)] $q\nu^2+1\equiv0\pmod D$;
\item[(ii)] $\dfrac{x}{cq}\equiv\dfrac{\nu}{D}+\dfrac{d}{cqD}\pmod 1$.
\item[(iii)] (Reciprocity.)  If $(k,m)=1$, $kx'\equiv1$ modulo $m$ and
$my'\equiv1$ modulo $k$, then
$\dfrac{x'}{m}+\dfrac{y'}{k}\equiv\dfrac1{km}\pmod1$.
\end{itemize}
\end{lemma}
\begin{proof}
(i) $q\nu^2+1=De^2-(de-1)(de+1)$.

(ii) The integer $xD-c^2qe-d$ equals $cq(cx-ce)+d(dx-1)$, so it is divisible
by $cq$.  It also equals $D(x-e)+d(de-1)$, so it is divisible by $D$.  Since
$(cq,D)=(cq,d^2)=1$, it is divisible by $cqD$.

(iii) The integer $kx'+my'-1$ is divisible by $m$ and by $k$, which are
coprime.
\end{proof}

In~\eqref{eq:G}, with $c=r\asymp L$, $d\asymp L^2$ and $q\asymp L^2$, one has
$D\asymp L^4$ and the correction in (ii) is of size $L^{-5}$.  Since
$n\le L^{3/2}$,
\[
 G_q(n)=\sum_{c\asymp L}\ \sum_{d\asymp L^2}
 e\Bigl(\frac{2n\,\nu(c,d)}{qc^2+d^2}\Bigr)\bigl(1+O(L^{-7/2})\bigr).
\]
So $G_q(n)$ is a Weyl sum over those roots of $q\nu^2+1\equiv0$ modulo $D$
that come from representations of $D\asymp q^2$ by the principal form
$qc^2+d^2$, at frequencies $n$ up to $D^{3/8}$.  This reflects the fact that
the vector $a$ of (F6) is essentially evaluation at the Heegner point
$i/\sqrt q$.  For a fixed polynomial, the equidistribution of such roots is
classical (Hooley; Duke, Friedlander and Iwaniec~\cite{DFI95};
T\'oth~\cite{Toth}).  (LS) is a mean-square bound for these Weyl sums that is
uniform in the coefficient $q$ up to $D^{1/2}$.

\medskip
\noindent\textbf{(F9) Moments of the zeta function.}  (F3) is one dyadic
piece of the off-diagonal term of the fourth moment of $\zeta$ twisted by a
Dirichlet polynomial with smooth coefficients of length $T^{\vartheta}$,
with $\vartheta=\theta/(2-\theta)$:
\[
 \theta=\tfrac12\leftrightarrow\vartheta=\tfrac13,\qquad
 \tfrac47\leftrightarrow\tfrac25,\qquad
 \tfrac35\leftrightarrow\tfrac37,\qquad
 \tfrac23\leftrightarrow\tfrac12 .
\]
The last is the sixth moment.  The integral in question is
\[
 \int_T^{2T}|\zeta(\tfrac12+it)|^4\Bigl|\sum_{m\asymp M}m^{-1/2-it}\Bigr|^2dt,
 \qquad M=T^{2/5+\delta},
\]
strictly between the fourth and sixth moments.  The correspondence is one
of sizes only.  (T) is the piece of this integral in which each divisor
function is restricted to factorizations of shape $(T/M)\cdot M$.  The
integral also contains the balanced factorizations
$\sqrt T\cdot\sqrt T$, which give the box $(2/7,5/14,5/14)$; that box is
closer to the cube and is not covered by the method.  So (T) does not
amount to an asymptotic formula for the integral.  For arbitrary
coefficients the record is $\vartheta<1/4$ (Watt~\cite{Watt};
Bettin, Bui, Li and Radziwi\l\l~\cite{BBLR}).

\medskip
\noindent\textbf{Other descriptions.}
\begin{itemize}
\item \emph{A determinant.}  $b\ell_1\ell_2-an_1n_2$ is the determinant of
$\left(\begin{smallmatrix}b\ell_1&an_1\\ n_2&\ell_2\end{smallmatrix}\right)$,
and $xyz+uvw$ is the determinant of the $3\times3$ matrix with rows
$(x,u,0)$, $(0,y,v)$, $(w,0,z)$.  So (T) counts integer matrices of
determinant $2$ in $\mathrm{GL}_3$ with three prescribed zero entries.
\item \emph{Farey fractions.}  (T) counts pairs of fractions
$b/(n_1n_2)$, $a/(\ell_1\ell_2)$ in a short arc that are adjacent in the
sense of~\eqref{eq:T}, weighted by restricted divisor functions of the two
denominators.
\item \emph{Almost-primes.}  After every piece with a flexible factor is
removed, what remains is the number of $n\le N$ such that $n$ and $n+2$ each
have two prime factors in $[P,2P]$, $P\approx N^{2/7}$.  All lower mixed
moments of the two counts are known.
\end{itemize}

\subsection{Why $4/7$}\label{sec:why47}
Three methods of different nature stop at the same exponent.  Write the
modulus as $q=rs$.
\begin{itemize}
\item \emph{Deshouillers--Iwaniec.}  The binding conditions are $H\le s$,
$H\le\sqrt r$ and $H\le r/\sqrt s$.
\item \emph{Polymath8a, Type~III}~\cite[\S8]{Polymath8a}.  Complete all three
variables to obtain $\mathrm{Kl}_3(2h_1h_2h_3;q)$, apply Cauchy--Schwarz in
$h$, and use Deligne's bounds for the correlation sums.  The factor $r$ must
satisfy $N^{4\theta-2}\le r\le N^{2-3\theta}$.  We compared with the model
computation of that paper, not with its final theorem.
\item \emph{Bessel's inequality in} (F6).
\end{itemize}

\begin{lemma}\label{lem:tight}
\leavevmode
\begin{itemize}
\item[(i)] At $r=s=L$ the three conditions of Deshouillers and Iwaniec all
hold if and only if $H\le q^{1/4}$, that is, $\theta\le4/7$.
\item[(ii)] The interval $[4\theta-2,\,2-3\theta]$ is nonempty if and only
if $\theta\le4/7$.  At $\theta=4/7$ it is the single point $2/7=\theta/2$,
the balanced factorization.
\end{itemize}
\end{lemma}
\begin{proof}
Linear arithmetic in $\theta$.
\end{proof}

The second method is purely algebraic.  It uses neither the shape of the
box nor the smooth averaging over $r$ and $s$.  The first discards the
divisor structure of $d=n_1n_2$ and treats each level $ss'$ separately.
The bound of Fouvry, Kowalski and Michel for $\sum_nd(n)K(n)$ uses only
$(n_1,n_2)$ and saves $q^{1/16}$ for a single frequency and a prime modulus.
We found no arrangement that combines two of these.

\begin{claim}[A no-go count]\label{claim:nogo}
Consider (F2) with $h$ of size $q^{\kappa}$ and $u,v,r,s$ of size
$q^{1/2}$.  Allow the following steps: Cauchy--Schwarz with an inner set of
variables $I$, then completion of the product of the outer dual variables to
the modulus $M$ of the resulting correlation sum, with square-root
cancellation assumed for every complete sum.  Then the saving is at most
$q^{1/4}$.  A second Cauchy--Schwarz does not improve this.
\end{claim}

\emph{Derivation.}  In this model
$\text{saving}^2=\min(|I|,\sigma)$ with
$\sigma=\min(\sqrt M,N_O/\sqrt M)$, where $N_O\le q^{5/4}$ is the length of
the outer dual product and $M$ is $q$ times the product of the inner
modulus variables.  A saving beyond $q^{1/4}$ needs $|I|>q^{1/2}$ and
$\sigma>q^{1/2}$.  The second forces $M<N_O^2/q\le q^{3/2}$, so $I$ contains
at most one of $r,s$.  Then $|I|>q^{1/2}$ forces a dual variable into $I$,
which gives $N_O\le q$ and $\sigma\le q/\sqrt M\le q^{1/4}$.  After a
second Cauchy--Schwarz on the off-diagonal, $\sigma$ is replaced by
$\sqrt{\min(|I_2|,\sigma_2)}$ with $\sigma_2\le\sqrt{N_O}\le q^{5/8}$, which
is below $q^{1/2}$.

\medskip
So a proof past $4/7$ needs a step outside this list.  The following were
examined and do not help.
\begin{itemize}
\item \emph{Shifting a smooth variable} (Burgess, $q$-van der Corput).  In
(F1) and (F2) no smooth variable is longer than a factor of the modulus it
meets.  In (F5) the frequency $n$ is, once the sum is extended to all
$n\le L^{3/2}$.  One differencing step then saves $L^{1/4}$ on the terms
of $S(q)$ with equal modulus variables $r=r'$, where $L^{1/2+\delta}$ is
needed, and nothing on the terms with $r\ne r'$.
\item \emph{Splitting each smooth variable} into a part with a divisor in
the useful range and the rest.  The rest is a prime times a small number;
expanding the prime by any exact identity returns the same smooth piece
with coefficient~$1$, so the recursion does not contract.
\item \emph{A variance over the long variable} $a$: the cross term is the
original sum.
\item \emph{Dirichlet characters}: a factor $\sqrt H$ is lost against the
trivial bound.
\item \emph{Positivity}: the known mixed moments of the two almost-prime
counts do not determine the missing one.
\item \emph{A divisor weight on the modulus through oldforms}: the family
becomes all newforms of conductor up to $q$, of the same size as before.
\item \emph{$\mathrm{GL}_3$.}  Sums of $\mathrm{Kl}_3$ over the modulus are
one term of the Kuznetsov formula for $\mathrm{SL}_3(\Z)$, isolated by
Buttcane~\cite{But}.  By the count given there, the trivial bound on the
spectral side is worse than Deligne's bound on the Kloosterman side.  Also
the modulus there is one smooth variable, which corresponds to
$d_3\times d_2$; ours is $\ell_1\ell_2$.  The long-element Kloosterman sums
factor into classical sums with first-power cross inverses, whereas ours
have squared inverses.
\item \emph{A third expression for the spectral sum in} (F6).  Poisson
summation, reciprocity and a second application of Kuznetsov's formula each
return the same sum, up to the Atkin--Lehner sign: at $\theta=4/7$ both
sides of the formula have $\asymp q$ terms.
\end{itemize}

\subsection{Numerical experiments}\label{sec:numerics}
\emph{Experiment 1: the variance in} (F4) (\texttt{notes/lab/var\_nu.c}).
For $600$ random odd $d\asymp L^2$ and all odd $r\asymp L$ coprime to $d$,
the program computes $\nu(d,r)$ of~\eqref{eq:nu} with $a\asymp L^{\alpha}$.
Here $\alpha=2(1-\theta)/\theta$, so $\alpha=1.5$ is $\theta=4/7$ and
smaller $\alpha$ is past the barrier.

\medskip
\begin{center}
\begin{tabular}{rrrrrrr}
$\alpha$ & $\theta$ & $L$ & $K$ & mean of $\nu$ & variance & variance/mean\\ \hline
$1.5$ & $0.571$ & $32$ & $5.7$ & $1.15$ & $1.14$ & $1.02$\\
$1.5$ & $0.571$ & $64$ & $8.0$ & $1.65$ & $1.68$ & $1.02$\\
$1.5$ & $0.571$ & $128$ & $11.3$ & $2.31$ & $2.46$ & $1.07$\\
$1.5$ & $0.571$ & $256$ & $16.0$ & $3.23$ & $3.52$ & $1.09$\\
$1.4$ & $0.588$ & $256$ & $9.2$ & $1.85$ & $2.15$ & $1.15$\\
$1.3$ & $0.606$ & $256$ & $5.3$ & $1.07$ & $1.29$ & $1.20$
\end{tabular}
\end{center}
\medskip

The variance is of the size of the mean, as for a Poisson variable, on both
sides of the barrier.  The target is $K^{2-\delta}$.

\emph{Experiment 2: the mean square in} (F5) (\texttt{notes/lab/rls.c}).
For random coprime $s,s'\asymp L$ the program computes $G_q(n)$
of~\eqref{eq:G} for $1\le n\le4HL$, with $P$ the number of terms
in~\eqref{eq:G}.  Square-root cancellation is $|G_q(n)|^2\approx P$.  The
entries are means of $|G_q(n)|^2/P$.

\medskip
\begin{center}
\begin{tabular}{rrrrrr}
$L$ & pairs & all $n$ & $(n,ss')=1$ & $n\in\mathcal N(s,s')$ &
$n\in\mathcal N$, $(n,ss')=1$\\ \hline
$16$ & $40$ & $0.98$ & $1.19$ & $0.99$ & $1.16$\\
$32$ & $40$ & $0.96$ & $1.11$ & $0.91$ & $1.02$\\
$64$ & $40$ & $0.99$ & $1.15$ & $1.03$ & $1.22$\\
$128$ & $16$ & $0.99$ & $1.14$ & $1.03$ & $1.15$\\
$256$ & $3$ & $0.98$ & $1.09$ & $0.99$ & $1.09$
\end{tabular}
\end{center}
\medskip

So $G_q(n)$ has square-root size for essentially every $n\ne0$ in the range,
on the Bohr set and off it, and $S(q)\approx L^{9/2}$.  The signed sums over
the Bohr set are also of square-root size.  Both targets are true with a
wide margin, in the range computed.

\emph{Experiment 3: the level, and the two parts of $S(q)$}
(\texttt{notes/lab/levelscan.c}).  Write $G_q=\sum_rG_r$ with
$G_r(n)=\sum_de(2n\inv d/(rq))$.  Then
$S(q)=\sum_r\|G_r\|^2+\sum_{r\ne r'}\langle G_r,G_{r'}\rangle$, a diagonal
part and a cross part.  The entries are divided by the number of
frequencies times $P$, so square-root cancellation is $1$.

\medskip
\begin{center}
\begin{tabular}{rrrrr}
$L$ & $q$ & total & $r=r'$ & $r\ne r'$\\ \hline
$64$ & $1$ & $25.82$ & $51.62$ & $-25.80$\\
$64$ & $7$ & $6.92$ & $6.90$ & $0.03$\\
$64$ & $65$ & $1.17$ & $0.66$ & $0.51$\\
$64$ & $513$ & $0.88$ & $0.90$ & $-0.03$\\
$64$ & $4097$ & $1.09$ & $0.99$ & $0.10$\\
$64$ & $8191$ & $1.00$ & $1.00$ & $0.01$
\end{tabular}
\end{center}
\medskip

For $q\gtrsim L$ the sum $S(q)$ is essentially its trivial diagonal, the
terms with the same $r$ and the same $d$.  At level $q=1$ there is no
square-root cancellation, and the cross part is $-\tfrac12$ of the diagonal
part.  So the diagonal part is not bounded by $S(q)$, and (LS) contains two
separate problems: a bound $L^{5-\delta}$ for the diagonal part, and the
same bound for the cross part.
\begin{itemize}
\item The diagonal part is a statement about one modulus $M=rq$:
$\sum_{n\le M^{1/2}}|\sum_{d\asymp M^{2/3}}e(n\inv d/M)|^2\ll M^{4/3-\delta}$
on average.  Weil's bound gives $M^{3/2}$, one van der Corput step
$M^{17/12}$, and the observed size is $M^{7/6}$.  Equivalently, the number
of $(d,d',j)$ with $\inv d-\inv{d'}\equiv j$ modulo $M$ and
$|j|\le M^{1/2}$ is asymptotic to its expected value with a power saving.
The spectral method says nothing about this part alone.
\item The cross part is where the cancellation over $r$ given by
Kuznetsov's formula is indispensable.
\end{itemize}

\emph{Experiment 4: the variance at all scales}
(\texttt{notes/lab/arcs.c}).  Let $N(x,\ell)$ be the number of points
$2\inv d/(rq)$ of~\eqref{eq:G} in the arc $[x,x+\ell)$.  By Parseval,
$\operatorname{Var}_xN(x,\ell)\approx\ell^2\sum_{0<|n|\lesssim1/\ell}|G_q(n)|^2$,
so the dual large sieve gives
\[
 \operatorname{Var}_xN(x,K/q)\ll q^{\eps}K\cdot\text{mean}\qquad(1\le K\le q).
\]
At $K=1$ this is a variance of Poisson order, and it is sharp.  For larger
$K$ it is what the case $K=1$ gives if the counts in the $K$ sub-arcs are
fully correlated.  (LS) is the statement
$\operatorname{Var}\ll K^{1-\delta}\cdot\text{mean}$ at $K=q^{1/4}$.  So the
barrier at $4/7$ is the absence of any information on the correlation
between neighbouring arcs of length $1/q$.  The table gives the ratio of
variance to mean for arcs of length $K/q$, and the largest correlation
coefficient between arcs at lags $1$ to $8$.

\medskip
\begin{center}
\begin{tabular}{rrrrrrrr}
$L$ & $q$ & $K=1$ & $K=4$ & $K=16$ & $K=64$ & $K=256$ & lags $1$--$8$\\ \hline
$64$ & $4097$ & $1.07$ & $1.10$ & $1.06$ & $1.10$ & $1.06$ & $0.025$\\
$64$ & $8191$ & $1.02$ & $0.98$ & $1.00$ & $0.85$ & $0.94$ & $0.015$\\
$128$ & $16385$ & $0.97$ & $0.93$ & $0.88$ & $1.01$ & $1.18$ & $0.013$
\end{tabular}
\end{center}
\medskip

The variance equals the mean at every scale, and neighbouring arcs are
uncorrelated to within the noise.

\emph{Experiment 5: the Kloosterman block as an operator}
(\texttt{notes/lab/kscan.c}, \texttt{kscan2.c}, \texttt{krow.c}).  Poisson
summation over $d$ in~\eqref{eq:G} gives
$G_q(n)=\sum_m\sum_cW(m,c)\,S(m,2n;cq)/(cq)$ with a smooth weight $W$
supported on $|m|\lesssim L$, $c\asymp L$; the term $m=0$ is a sum of
Ramanujan sums.  The experiment concerns the kernel with the weight removed,
\begin{equation}\label{eq:block}
 K(m,n)=\sum_{c\asymp\sqrt q}\frac{S(m,n;cq)}{cq},\qquad
 m\asymp\sqrt q,\quad 1\le n\le N=q^{3/4}.
\end{equation}
If $\Gamma_0(q)$ has no exceptional eigenvalues, Kuznetsov's formula and
Bessel's inequality for the forms of level $q$ (the large sieve
of~\cite{DI82}, both ranges being below the level) give
$\|KK^{T}\|\ll q^{\eps}$; the bound $S(q)\ll L^{5+\eps}$ is the corresponding
statement for one smooth vector, and for that vector it is elementary.  In
mean square the entries of $K$ have the size $q^{-1/2}$, so the mean
eigenvalue of $KK^{T}$ is $\asymp N/q$, smaller than the bound for the norm
by the factor $q/N=q^{1/4}$.  (LS) says that the smooth vector is not near
the top of the spectrum.  The table gives the
largest eigenvalue of $KK^{T}$ divided by the mean eigenvalue, for all prime
levels in four ranges ($c,m\in[\sqrt q,2\sqrt q)$, sharp cut-offs).

\medskip
\begin{center}
\begin{tabular}{lrrrr}
prime levels near & $280$ & $1000$ & $4100$ & $16000$\\ \hline
number of levels & $16$ & $30$ & $51$ & $29$\\
median of top/mean & $2.78$ & $3.09$ & $3.50$ & $3.24$\\
maximum of top/mean & $4.14$ & $5.03$ & $4.68$ & $3.72$\\
edge for independent entries, $(1+\sqrt{M/N})^2$ & $2.23$ & $2.00$ & $1.83$ & $1.69$\\
Bessel, $1+q^{1/4}$ & $5.1$ & $6.6$ & $9.0$ & $12.2$
\end{tabular}
\end{center}
\medskip

The spectrum is that of a matrix with independent entries, with one or two
eigenvalues above the edge, and the largest eigenvalue does not grow with
$q$.  Composite levels behave in the same way (products of two primes near
$\sqrt q$: $2.1$ to $3.7$; squares: $3.2$ to $3.5$; smooth levels: $2.2$ to
$2.5$, for $q$ between $960$ and $1770$).  For single rows,
$\sum_{n\le N}K(m,n)^2$ divided by its generic value
$N\sum_c\varphi(cq)/(cq)^2$ lies between $0.88$ and $1.17$ for
$m=1,\dots,8$ at $q=1009$, $4001$, $4003$, $16001$, $64007$, where the ratio
of Bessel's bound to the generic value is $5.6$, $8$, $8$, $11.3$, $15.9$.
So in the range computed the missing estimate holds for every vector, and
not only for the smooth one.  We record the corresponding statements.

\medskip\noindent\textbf{(OP).}  \emph{There is $\delta>0$ such that for
all complex $a_m$ and $b_n$, and $g$ smooth and supported on
$c\asymp\sqrt q$,
\[
 \sum_{m\asymp\sqrt q}\ \sum_{n\le q^{3/4}}a_mb_n\sum_{c}g(c)
 \frac{S(m,n;cq)}{cq}\ll q^{-\delta}\|a\|\,\|b\|.
\]}

\noindent\textbf{(OP$_1$).}  \emph{There is $\delta>0$ such that
$\displaystyle\sum_{n\le q^{3/4}}\Bigl|\sum_cg(c)\frac{S(1,n;cq)}{cq}
\Bigr|^2\ll q^{-\delta}$.}

\medskip
For (OP) the method of~\cite{DI82} gives $q^{\eps}\|a\|\|b\|$ if there are
no exceptional eigenvalues, and a larger power of $q$ otherwise; Weil's
bound gives $q^{3/8}\|a\|\|b\|$; and a bound $q^{-1/2+\eps}$ for every entry,
which is the conjecture of Linnik and Selberg in the level aspect, gives
only $q^{1/8+\eps}\|a\|\|b\|$.  The expected size is $q^{-1/8}$.  The
analogue of (OP) for each dyadic range $m\asymp M\le\sqrt q$ would give
(LS) once the variables $m$ and $c$ are separated in the weight $W$; we
have not written out the details.  (OP$_1$) is the case of a single row:
Bessel's inequality gives $O(1)$ on the same hypothesis, Weil's bound
$q^{1/4}$, and the expected size is $q^{-1/4}$.  It is the smallest instance of the barrier
that we know.  Its off-diagonal part has the following closed form.

\begin{lemma}\label{lem:corr}
Let $q$ be prime, $c,c'\ge1$ with $(c,c')=(cc',q)=1$, $k=cc'q$,
$(h,k)=1$, and let $m,m'$ be integers prime to $q$.  Then
\[
 \sum_{n\bmod k}S(m,n;cq)\,S(m',n;c'q)\,\e\Bigl(\frac{hn}{k}\Bigr)
 =k\,\Bigl[\e\Bigl(-\frac{A\,\inv h}{k}\Bigr)S(1,mm'\inv h^{\,2};q)
 -\e\Bigl(-\frac{A\,\inv{hq}}{cc'}\Bigr)\Bigr],\qquad A=mc'^2+m'c^2,
\]
where $\inv h$ is an inverse of $h$ modulo $k$ and $\inv{hq}$ an inverse
of $hq$ modulo $cc'$.
\end{lemma}

\begin{proof}
Opening both sums, the left-hand side is $k$ times the sum of
$\e(mx/(cq)+m'y/(c'q))$ over reduced $x\bmod cq$, $y\bmod c'q$ with
$c'\inv x+c\inv y+h\equiv0\pmod k$.  Modulo $c$ and $c'$ the condition
determines $x\equiv-c'\inv h$ and $y\equiv-c\inv h$, which gives the factor
$\e(-A\,\inv{hq}/(cc'))$.  Modulo $q$, put $X=x\inv c$, $Y=y\inv{c'}$
and $t=hcc'X+c'^2$; the condition is $Y=-c\inv{hc'}+cc'\inv h\,\inv t$, and
$X\not\equiv0$, $Y\not\equiv0$ hold exactly when
$t\not\equiv0,c'^2$.  Since $mX+m'Y=m\inv{hcc'}\,t+m'cc'\inv h\,\inv t
-A\,\inv{hcc'}$, the sum over $t\not\equiv0$ is
$\e(-A\,\inv{hcc'}/q)\,S(1,mm'\inv h^{\,2};q)$, and the term
$t\equiv c'^2$, which must be removed, equals $1$.  The two factors
combine by the Chinese remainder theorem.
\end{proof}

The identity was also checked numerically (to $10^{-9}$, for eleven choices
of $(c,c',q,m,m')$ and all $h$).  It is notable that the Kloosterman sum on
the right does not depend on $c$ and $c'$.  By Poisson summation, the part
of the sum in (OP$_1$) (the case $m=m'=1$) with $c\ne c'$ coprime is
therefore, up to a term of lower order and the terms with $(h,k)>1$,
\[
 \frac N{q^2}\sum_{h\ll q^{5/4}}S(1,\inv h^{\,2};q)\sum_{c\ne c'}
 \frac{g(c)g(c')}{cc'}\,\widehat w\Bigl(\frac{hN}{cc'q}\Bigr)\,
 \e\Bigl(\inv q\,\frac{c\inv{c'}+\inv cc'}{h}\Bigr)
 \e\Bigl(-\frac{c^2+c'^2}{hcc'q}\Bigr),
\]
with inverses modulo $h$ in the last exponential but one.  The trivial bound
for this is $q^{1/2}$, Bessel's inequality corresponds to $O(1)$, and
(OP$_1$) to $q^{-1/4-\delta}$.  Since
$S(\inv q,\inv q;h)\,S(\inv h,\inv h;q)=S(1,1;hq)$, the sum over $c,c'$
and the Kloosterman sum combine into the sum $S(1,1;hq)$ with its variable
of summation restricted, modulo $h$, to the ratios $c/c'$ of two integers of
size $\sqrt q$; these are the residues $\rho$ for which $\rho/h$ lies within
$O(1/h)$ of a fraction with denominator of size $\sqrt q$, a set of density
about $q^{-1/4}$.

\emph{Experiment 6: the block modulus by modulus}
(\texttt{notes/lab/kmod.c}).  Write $K=\sum_cK_c$ with
$K_c(m,n)=S(m,n;cq)/(cq)$.  The table gives largest eigenvalues divided by
the mean eigenvalue of $KK^{T}$ (in the last column, of $K_cK_c^{T}$).

\medskip
\begin{center}
\begin{tabular}{rrrrr}
$q$ & $KK^{T}$ & $\sum_cK_cK_c^{T}$ & $\sum_c\|K_c\|^2$ &
$\max_c$ for $K_cK_c^{T}$\\ \hline
$1009$ & $3.36$ & $1.11$ & $1.82$ & $2.15$\\
$4001$ & $3.46$ & $1.15$ & $1.88$ & $2.01$\\
$16001$ & $2.58$ & $1.07$ & $1.69$ & $1.97$
\end{tabular}
\end{center}
\medskip

Each single modulus has the spectrum of a matrix with independent entries,
with no eigenvalue above the edge; the part of $KK^{T}$ that is diagonal in
$c$ is close to a multiple of the identity; and the eigenvalues above the
edge come from the terms $c\ne c'$, which Lemma~\ref{lem:corr} evaluates.
So (OP) is the sum of two statements, neither of which is known: a bound
$\|(S(m,n;k))_{m\le k^{1/3},\,n\le k^{1/2}}\|\ll k^{5/6-\delta}$ for the
single modulus $k=cq$ (Weil's bound gives $k^{11/12}$, the expected size is
$k^{3/4}$, and a fourth moment of singular values with an exact count would
give $k^{5/6}$), and the bound $q^{-\delta}$ for the norm of the terms
$c\ne c'$.

\emph{The same statement in $\mathrm{SL}_2(\Z/q\Z)$.}  Let a matrix
$\gamma=\bigl(\begin{smallmatrix}a&b\\c&d\end{smallmatrix}\bigr)$ act on
functions on $\Z/q\Z$ through $x\mapsto(ax+b)/(cx+d)$, and let
$\psi_m(x)=\e(mx/q)$.  For $q\nmid c$,
\[
 \langle\psi_n,\rho(\gamma)\psi_m\rangle
 =\frac1q\,\e\Bigl(\frac{(ma+nd)\inv c}{q}\Bigr)S(mn\inv c^{\,2},1;q).
\]
Summing over the $\varphi(c)$ matrices $\gamma_{c,d}\in\mathrm{SL}_2(\Z)$
with lower left entry $c$ and $0<a,d<c$, and using
$\inv c/q+\inv q/c\equiv1/(cq)$, one finds
$\sum_d\langle\psi_n,\rho(\gamma_{c,d})\psi_m\rangle
=q^{-1}S(m,n;cq)\,(1+O((|m|+|n|)/q))$.  So, up to this factor,
\[
 K=P_{q^{3/4}}\;\rho\Bigl(\sum_{c\asymp\sqrt q}\frac1c\sum_d\gamma_{c,d}\Bigr)
 \;P_{\sqrt q},
\]
where $P_V$ is the projection onto the span of the $\psi_m$ with
$|m|\le V$, which is an average over the unipotent elements
$\bigl(\begin{smallmatrix}1&h\\0&1\end{smallmatrix}\bigr)$ with
$|h|\le q/V$.  The sum in the middle runs over the $\asymp q$ elements of
$\mathrm{SL}_2(\Z)$ of norm $\asymp\sqrt q$, with total weight
$\asymp\sqrt q$.  The triangle inequality gives $\|K\|\ll\sqrt q$; the
bound $\|K\|\ll q^{\eps}$ is the statement that this ball of
$\mathrm{SL}_2(\Z)$ acts on the nontrivial part of $L^2(\mathbb
P^1(\Z/q\Z))$ with square-root cancellation, which is where the spectral
gap enters; and (OP) asks for a further saving from the two band limits.
In~\cite{Pas25b} the ball is replaced by the single element
$\bigl(\begin{smallmatrix}0&-1\\1&0\end{smallmatrix}\bigr)$, the saving
comes from the band limits alone, and the proof counts closed words in the
group; that count is elementary because the words do not fill the group.
Here a closed word of the $2k$-th moment has $q^{11k/4}$ choices against
$|\mathrm{SL}_2(\Z/q\Z)|\asymp q^3$, so that only $k=1$, the
Hilbert--Schmidt norm $q^{1/8}$, is in that regime.

\emph{Experiment 7: a model without the ball}
(\texttt{notes/lab/kmodel.c}).  By twisted multiplicativity
$S(m,n;cq)=S(m\inv q,n\inv q;c)\,S(1,mn\inv c^{\,2};q)$.  Dropping the first
factor leaves, for $q$ prime and $\mathrm{Kl}(t)=q^{-1/2}S(1,t;q)$,
\[
 K_0(m,n)=\sum_{c\asymp\sqrt q}\mathrm{Kl}(mn\inv c^{\,2}),\qquad
 m\asymp\sqrt q,\quad n\le N=q^{3/4}.
\]
Here nothing is summed over $\mathrm{SL}_2(\Z)$: $K_0$ is the restriction to
two intervals of a convolution operator on the cyclic group
$(\Z/q\Z)^{\times}$, whose multiplier at a character $\chi$ is
$q^{-1/2}\tau(\chi)^2\sum_c\chi^2(c)$.  Completing the variable $m\inv c^{\,2}$
and using the orthogonality of $u\mapsto\mathrm{Kl}(un)$ gives
$\|K_0\|\ll q^{3/4+\eps}$, which is the analogue of Bessel's inequality; the
expected size is $q^{5/8}$, apart from the quadratic character $\eta$, the
one character for which $\sum_c\chi^2(c)$ does not cancel, which contributes
$\pm\eta(m)\eta(n)$.  The table gives the three largest eigenvalues of
$K_0K_0^{T}$ over the mean eigenvalue, before and after removing that term.

\medskip
\begin{center}
\begin{tabular}{rrrrrr}
$q$ & before & coefficient of $\eta\otimes\eta$ & after & edge & $q/N$\\ \hline
$1009$ & $2.26,\ 1.74,\ 1.68$ & $0.973$ & $1.82,\ 1.74,\ 1.69$ & $2.01$ & $5.6$\\
$4001$ & $2.18,\ 1.79,\ 1.72$ & $0.943$ & $1.83,\ 1.77,\ 1.72$ & $1.83$ & $8.0$\\
$16001$ & $2.28,\ 1.70,\ 1.63$ & $1.054$ & $1.72,\ 1.64,\ 1.60$ & $1.68$ & $11.3$\\
$64007$ & $2.16,\ 1.60,\ 1.58$ & $-1.011$ & $1.61,\ 1.59,\ 1.56$ & $1.56$ & $15.9$
\end{tabular}
\end{center}
\medskip

So the model has the spectrum of a matrix with independent entries, to
within a few per cent, and the bound known for it is larger by the same
factor $q/N=q^{1/4}$ as for $K$.

\medskip\noindent\textbf{(MP).}  \emph{For $q$ prime and all $a$, $b$,
\[
 \sum_{m\asymp\sqrt q}\ \sum_{n\le q^{3/4}}\ \sum_{c\asymp\sqrt q}
 a_mb_n\,\mathrm{Kl}(mn\inv c^{\,2})
 =\pm\Bigl(\sum_ma_m\eta(m)\Bigr)\Bigl(\sum_nb_n\eta(n)\Bigr)
 +O\bigl(q^{3/4-\delta}\|a\|\,\|b\|\bigr).
\]}

\medskip
We could not prove (MP).  H\"older's inequality on the side of the
characters, with the fourth moments of the two character sums of length
$\sqrt q$ (which are known exactly), gives $q^{3/4}$ again; a shift
$c\mapsto c+rs$ followed by H\"older's inequality and complete sums of
products of Kloosterman sums gives $q$; and writing the excess as a sum over
$q^{5/4}$ twists of the square matrix
$(\mathrm{Kl}(mm'\xi))_{m,m'\asymp\sqrt q}$ of~\cite{KMS,Pas25b}, even the
best possible bound for each twist loses $q^{1/4}$ to (MP).  The conclusion
we draw is that the obstruction at $4/7$ is not the need to combine the
spectral gap with a property of finite groups: the model has no spectral
gap in it and is stuck at the same place.

For a smooth vector $a$ the model takes a concrete form.  The functional
equation of the smooth character sum over $m$ absorbs one Gauss sum, the
row sums become the Fourier transform of the $q$ points
$\inv{xc^2}\in\Z/q\Z$ with $x,c\asymp\sqrt q$, and the analogue of (LS) is
the statement that, on average over $0<|t|\le q^{1/4}$,
\[
 \#\bigl\{(x,c,x',c'):\ (1-txc^2)(1+tx'c'^2)\equiv1\pmod q\bigr\}
 =q\,\bigl(1+O(q^{-\delta})\bigr),
\]
all four variables being of size $\sqrt q$.  By positivity the left-hand
side is at most its value at $t=0$, the number of solutions of
$xc^2\equiv x'c'^2$, which is $\asymp q$ in the range computed; this is the
analogue of Bessel's inequality and is all that we can prove.  In that
range (\texttt{notes/lab/mcorr.c}, $q$ up to $10^7$) the left-hand side is
$q$ to within the statistical noise for every $t\ne0$.  So the missing estimate is any
power saving in an asymptotic count for one polynomial congruence in four
variables in a box of side exactly $\sqrt q$.

\subsection{Estimates examined that do not apply}\label{sec:notapply}
\begin{itemize}
\item \emph{Watt's mean value theorem}~\cite{Watt}, in the form used
in~\cite{BBLR}: applied directly to (F1) it saves $H^{1/2}$; $H^{1+\delta}$
is needed.
\item \emph{The quadratic divisor theorem of}~\cite{BBLR}: its error term
requires the product of the two coefficient ranges to be below $N^{1/2}$.
The possible assignments give $N^{2-2\theta}$, $N^{\theta}$ and
$N^{1-\theta/2}$, which is admissible only for $\theta<1/2$ or
$\theta>3/4$.
\item \emph{Levels of distribution of $d_3$}~\cite{HB86,FKM3,Sha}: below
$4/7$, as listed in Section~\ref{sec:twosided}.  Aydemir and
Boran~\cite{AB} reach $8/11$ on average over residue classes to a prime
modulus; (T) has one class for each modulus.
\item \emph{Convolution estimates of Maynard and Pascadi}~\cite{May2,Pas25},
with every partition of $\{\ell_1,\ell_2\}$ into three groups, and Pascadi's
triple-convolution estimate, which needs $D\ge EN^{\eps}$.
\item \emph{Irving's Type~II bound}~\cite{Irv14}: $8/15$ at the balanced
point.
\item \emph{Grimmelt--Merikoski, determinant equations with a periodic
weight}~\cite{GM24}.  The method expands a Poincar\'e series over the orbits
of one congruence subgroup and applies Cauchy--Schwarz between a sum of
Fourier coefficients and a sum of point evaluations.  For one pair of our
moduli it gives the individual error $L^{7/4}$ of (F3).  Its gain needs many
orbits of one fixed group.  A weight is invariant only if it is a condition
modulo the level, so each of our two divisibility conditions forces its
modulus into the level, which is at least $L^2$ in every arrangement.
\item \emph{Grimmelt--Merikoski, roots of quadratic
congruences}~\cite{GM25}.  Their Type~I estimate treats roots of
$a\ell^2+h\equiv0$ modulo $k$, $k\equiv0$ modulo $d$, for a fixed $h$, on
average over $d$.  The sum in (F8) is the case $d=h=q$: the matrix
$\left(\begin{smallmatrix}(q\nu^2+1)/D&-q\nu\\-q\nu&qD\end{smallmatrix}\right)$
has determinant $q$.  The hypotheses hold, but the bound is for the whole
sum over $d\le L^2$ and exceeds the main term for a single $d$.  Their
saving comes from averaging over $d$ with $h$ fixed; (F8) has one $d$ for
each $h$ and an average over $h$.
\item \emph{Lau}~\cite{Lau} treats $d_k\times d_2$, not $d_3\times d_3$.
\item \emph{Bilinear forms with Kloosterman sums to composite moduli}
(Pascadi~\cite{Pas25b}; see also~\cite{KMS} for prime moduli).  These are
bounds of the same kind as (OP), for arbitrary coefficients and lengths up
to the square root of the modulus.  Corollary~1.4 of~\cite{Pas25b} is a
version summed over the moduli divisible by a level $q$ with a factorization
$q=de$, which is our situation with $d=s$, $e=s'$.  With moduli of size
$q^{3/2}$ it gives $q^{5/12+\eps}\|a\|\|b\|$ for the sum in (OP), against
$q^{\eps}\|a\|\|b\|$ from the large sieve: the bound is taken modulus by
modulus, and no bound of that form can reach (OP), since the best possible
one, $\sqrt{(M+N)cq}$ for each modulus $cq$, gives $q^{1/8}\|a\|\|b\|$.
What (OP) needs is such orthogonality together with the cancellation over
$c$ that Kuznetsov's formula supplies.  The method of~\cite{Pas25b} bounds a
sixth moment of singular values and uses the multiplicities of the
representations of $\mathrm{SL}_2(\Z/c\Z)$; whether it has an analogue on
the spectral side is open.
\end{itemize}

\subsection{Open questions}\label{sec:open}
\begin{enumerate}
\item Prove (LS), or any bound $S(q)\ll L^{5-\delta}$ on average over $q$.
\item Prove (MP), the model of Experiment~7, in which only finite fields
occur.
\item Let $\nu$ be the counting function of the points of the model on
$\Z/q\Z$ and $C(\lambda)=\sum_x\nu(x)\nu(x+\lambda)$.  By the
Cauchy--Schwarz inequality over all shifts, the analogue of (LS) follows from
$\sum_{\lambda\ne0}(C(\lambda)-q)^2\ll q^{9/4-\delta}$, a bound for the
additive energy of the point set, equivalently for the fourth moment of its
Fourier transform over all frequencies.  No short range occurs in this
statement.  Numerically the left-hand side is $\approx q^2$
(\texttt{notes/lab/e4.c}).  Can it be proved?
\item Prove (OP$_1$), or disprove (OP).  Is there a level $q$ for which the
largest eigenvalue of the block~\eqref{eq:block} exceeds the mean eigenvalue
by a power of $q$?
\item Does the argument of~\cite{GM25}, with symmetric matrices of
determinant $h$, admit an average over $h$ with $d=h$?
\item On average over $q$, the pair correlation behind $S(q)$ becomes a sum
of a restricted $d_3$ over the values $uu'+1$, with the side condition that
$u/j$ has a convergent with denominator of size $L$.  Can the side condition
be removed without recreating a level of size $L^2$?
\item Check Claims~\ref{claim:conrey}--\ref{claim:conditional} in detail.
The unconditional part of Claim~\ref{claim:conditional}, a level $3/5$ for
twin correlations of well-factorable weights, would be a theorem beyond the
bar of Claim~\ref{claim:conrey} for a restricted class of weights.
\item Formalize Kuznetsov's formula and the large sieve
of~\cite{DI82}, which would bring Claim~\ref{claim:conrey} within reach of
machine checking.
\end{enumerate}

\section{Formalization}\label{sec:lean}

The repository accompanying this paper contains a Lean~4 project
(\texttt{lean/}, Lean v4.35.0, Mathlib commit \texttt{56a00585}).  In it the
following are proved with no hypotheses:
\begin{itemize}
\item Theorem~\ref{thm:LambdaR}:
\texttt{MobiusDeterminant.lambdaR\_twin\_correlation\_unconditional};
\item Theorem~\ref{thm:main}, \eqref{eq:main-power} and the bound for
$\cS^0_\kappa$: \texttt{Determinant.detSum\_sub\_zeroMode},
\texttt{Determinant.zeroModeBound};
\item Theorem~\ref{thm:main}, \eqref{eq:main-log}:
\texttt{Determinant.detSum\_cancellation};
\item Corollary~\ref{cor:diagonal}: \texttt{Determinant.detSum\_diagonal}.
\end{itemize}
The command \texttt{\#print axioms} reports only Lean's standard axioms
(\texttt{propext}, \texttt{Classical.choice}, \texttt{Quot.sound}) for these
theorems and their entire dependency cones, including the external libraries
below.  No \texttt{sorry} appears anywhere.

\emph{Definitions.}  The file \texttt{Basic.lean} defines the objects of
Section~\ref{sec:LambdaR} directly: $\Lambda_R$, $\varrho(d,e)$, $r_N(d,e)$,
$M(R)$ and $S_g(y)$.  It defines the twin-prime constant as the Euler
product $\prod_p\mathfrak S_p$, with $\mathfrak S_2=2$ and
$\mathfrak S_p=1-(p-1)^{-2}$ for odd $p$.  The file
\texttt{Determinant/Defs.lean} defines $\cS_\kappa$ by~\eqref{eq:S}, with $W$
smooth and supported in $[1,2]^4$.  It also defines $\cS^0_\kappa$
by~\eqref{eq:zero-mode}, and $\Omega$ and $\sigma$ as in
Section~\ref{sec:intro}.  The statements are in
\texttt{Determinant/Statement.lean}.

\emph{Structure.}  The proof is split into components with fixed formal
interfaces (\texttt{Unconditional/Interfaces.lean}):

\medskip
\begin{center}
\begin{tabular}{ll}
Paper & Lean \\ \hline
Theorem~\ref{thm:LambdaR} & \texttt{lambdaR\_twin\_correlation\_unconditional}\\
Theorem~\ref{thm:main} & \texttt{detSum\_sub\_zeroMode}, \texttt{detSum\_cancellation}\\
Corollary~\ref{cor:diagonal} & \texttt{detSum\_diagonal}\\
Section~\ref{sec:proof}, product weights & \texttt{detTensorBound}, \texttt{oddMainBound}\\
Theorem~\ref{thm:level} ($\theta<9/17$) & \texttt{mobiusPairLevel}\\
Theorem~\ref{thm:meanvalue} & \texttt{meanValueThm}\\
Proposition~\ref{prop:bilinear} and Steps~1, 3 & \texttt{balancedBoxBound\_of\_BM}\\
Lemma~\ref{lem:BM} ($m=2$, loss $Q^\eps$) & \texttt{bloomMaynardE4}\\
Lemmas~\ref{lem:poisson}, \ref{lem:separation} & \texttt{poissonAPDecay}, \texttt{mellinSeparation}\\
\eqref{eq:PNTmu} & \texttt{pntMobius}\\
Section~\ref{sec:LambdaR} & \texttt{lambdaR\_twin\_correlation}\\
Lemmas~\ref{lem:HB}, \ref{lem:KC}, \ref{lem:energy-integers}, \ref{lem:energy-fractions} & \texttt{Core/}\\
all exponent computations of Sections~\ref{sec:intro}--\ref{sec:LambdaR} & \texttt{Exponents.lean}\\
Section~\ref{sec:barrier}: arithmetic and Lemmas & \texttt{Barrier.lean}
\end{tabular}
\end{center}
\medskip

\emph{Section~\ref{sec:barrier}.}  The file \texttt{Barrier.lean} checks the
parts of Section~\ref{sec:barrier} that are arithmetic or elementary.  It
does not formalize any of the analytic estimates quoted there, so each
Claim remains unverified.  A statement such as \texttt{conrey\_error\_iff}
says that an error term of the stated shape is a power saving exactly for
$\theta<4/7$; it does not say that the error term has that shape.

\medskip
\begin{center}
\begin{tabular}{ll}
Section~\ref{sec:barrier} & \texttt{Barrier.lean} \\ \hline
Lemma~\ref{lem:dictionary} and the table & \texttt{saving\_suffices\_iff}, \texttt{level\_table}\\
Remark~\ref{rem:comparison} (Irving's region) & \texttt{irving\_diagonal\_iff}, \texttt{irving\_crossing}\\
Claim~\ref{claim:conrey}, exponents only & \texttt{conrey\_error\_iff}\\
Claim~\ref{claim:window}, the window & \texttt{window\_nonempty\_iff}, \texttt{window'\_nonempty\_iff}\\
Lemma~\ref{lem:balanced-window} & \texttt{balanced\_mem\_window\_iff}, \texttt{balanced\_mem\_window'\_iff}\\
units at $\theta=4/7$ & \texttt{dual\_length\_eq\_iff}, \texttt{mean\_eq\_iff}\\
(F4), pairs of solutions & \texttt{pair\_of\_solutions}\\
Claim~\ref{claim:LS}, exponents only & \texttt{ls\_to\_pairs}, \texttt{pairs\_to\_dual}, \texttt{ls\_sizes}\\
(F5), pairs of fractions & \texttt{fraction\_pair}\\
Lemma~\ref{lem:roots} & \texttt{quadratic\_root}, \texttt{kloosterman\_fraction\_eq\_root}, \texttt{reciprocity}\\
(F9), the dictionary and the two boxes & \texttt{moment\_table}, \texttt{momentLength\_inv}, \texttt{moment\_boxes}\\
the two determinants & \texttt{det\_two}, \texttt{det\_three}\\
Lemma~\ref{lem:tight} & \texttt{di\_conditions\_balanced\_iff}, \texttt{polymath\_window\_nonempty\_iff}\\
Claim~\ref{claim:nogo}, the two inequalities & \texttt{nogo\_modulus\_bound}, \texttt{nogo\_saving\_bound}\\
Section~\ref{sec:why47}, van der Corput in the frequency & \texttt{vdc\_frequency}\\
Section~\ref{sec:notapply}, sizes & \texttt{bblr\_range\_iff}, \texttt{gm\_sizes}, \texttt{level\_loss\_iff}
\end{tabular}
\end{center}
\medskip

\emph{External formalizations used.}  The following are imported, not
re-proved:
\begin{itemize}
\item from Lean Pool: the Weil bound for Kloosterman sums to squarefree
moduli (adapted there from C.~Dorsey's Weil Bounds Toolkit), and the
divisor bound $\tau(n)\le n^{\eps}$;
\item from PNT+ (Kontorovich, Tao and contributors): the classical
zero-free region of $\zeta$ and bounds for $\zeta'/\zeta$ and $1/\zeta$;
\item also from PNT+: the smoothed Mellin--Perron pipeline, which is adapted
here from $-\zeta'/\zeta$ to $1/\zeta$.
\end{itemize}
To our knowledge, the formal proofs of the Bloom--Maynard energy bound and
of the prime number theorem for $\mu$ with error term
$x\exp(-c\sqrt{\log x})$ are new.

\emph{Differences from the written proof.}  None of these affects the
statements.
\begin{enumerate}
\item In the balanced boxes, the formal Korolev--Changa inequality is
applied to arbitrary residues $h\inv u$ modulo $e$, grouped by the reduced
fraction $h/u$.  This makes the gcd strata $g=(h,e)$ of Step~1 of
Theorem~\ref{thm:meanvalue} unnecessary.
\item In the long boxes, sharp dyadic intervals and completion over
$\Z/q\Z$, together with the Weil bound, replace the smooth Poisson step of
Lemma~\ref{lem:incomplete}.  The resulting $\log q$ loss is absorbed.
\item In Proposition~\ref{prop:smooth-level}, the coupled weight depends on
$(h,d',e')$ only through $h/(d'e')$.  One-variable Mellin separation
therefore suffices.
\item The formal prime number theorem for $\mu$ uses a bound
$1/\zeta(\sigma+it)\ll(1+|t|)^{1/2}$ in the zero-free region, obtained by
integrating $\zeta'/\zeta$ (Gr\"onwall).  This is weaker than the classical
$(\log|t|)^A$ but suffices.
\item For Theorem~\ref{thm:main}, the formal proof does not use Poisson
summation in $a$ modulo $2e$.  Instead it writes $n=ad$, so that
$\cS_\kappa$ becomes a smooth sum over odd $n$ with $d\mid n$ and
$e\mid n+\kappa$.  It then expands $W$ by a four-variable Mellin transform
(through Fourier inversion for Schwartz functions).  This gives a superposition of
product weights in $(n,d,e)$, and each one is treated by Poisson summation
in $n$ and Theorem~\ref{thm:meanvalue}, as in
Section~\ref{sec:level} but in the unbalanced range.  The exponents are those
of Section~\ref{sec:proof}.  The zero frequency~\eqref{eq:zero-mode} is
identified with the main term by the substitution $n/X=\alpha\,d/D$.  The
case $\kappa=-2$ is reduced exactly to $\kappa=2$ by the symmetry of
Section~\ref{sec:proof}.
\end{enumerate}

\emph{Trust.}  The Lean kernel guarantees that the formal theorems follow
from the formal definitions.  A human reader should still check that the
definitions in \texttt{Basic.lean} and \texttt{Determinant/Defs.lean} encode
the objects of this paper.  For Section~\ref{sec:barrier} the kernel checks
only the statements listed in the second table.

\section{Remarks}\label{sec:remarks}

\subsection{The boundary of the method}
The region $\Omega$ is exactly where the balanced estimate of
Theorem~\ref{thm:meanvalue} beats the Cauchy--Schwarz target.  The target is
$\cX\ll DE^{1/2}X^{-\delta}$ at $H\asymp T$, and the balanced bound meets it
precisely when $T<D^{1/9}$.  The bottleneck is case~1 of Step~3, at
$U\asymp D^{1/3}$, the shortest factor that the Heath-Brown decomposition
can force on the balanced side.  A finer treatment of factors near
$D^{1/3}$, or cancellation among the nonzero frequencies $h$ rather than
their bounded multiplicity, would be needed to go further.  Irving's Type~II
estimate~\cite{Irv14} is one such treatment.  It inverts the Kloosterman
fractions and applies Weil's bound to short Kloosterman sums.  For $E\asymp
D$ and $H<D^{1/3}$ it is stronger than Proposition~\ref{prop:bilinear} at
$U\asymp D^{1/3}$; see Remark~\ref{rem:comparison}.  The spectral method
appears to go further still, to $T<D^{1/4}$ (Claim~\ref{claim:conrey}), and
stops there for the reasons given in Section~\ref{sec:why47}.

\subsection{Another view: correlations along dilations}
Put $n=ad$, so that $n+\kappa=be$.  Then~\eqref{eq:S} becomes
\[
 \cS_\kappa=\sum_{a,b\ \mathrm{odd}}\
 \sum_{\substack{n\equiv0\ (\mathrm{mod}\ a)\\ n\equiv-\kappa\ (\mathrm{mod}\ b)}}
 \mu\Bigl(\frac na\Bigr)\mu\Bigl(\frac{n+\kappa}b\Bigr)
 W\Bigl(\frac aA,\frac bB,\frac n{aD},\frac{n+\kappa}{bE}\Bigr),
\]
where $n/a$ and $(n+\kappa)/b$ are also required to be odd.  This is a
two-point correlation of $\mu$ along the linear forms $n/a$ and
$(n+\kappa)/b$, averaged over dilations $a\asymp A$ and $b\asymp B$.  The
two extremes of the family are classical:
\begin{itemize}
\item when $x+y>1$, the Chinese remainder theorem suffices;
\item when $A=B=1$, the sum is essentially $\sum_n\mu(n)\mu(n+2)$, the open
binary case of Chowla's conjecture.  Only logarithmically averaged and
shift-averaged versions of that case are known~\cite{Tao16,MRT15}.
\end{itemize}
Theorem~\ref{thm:main} extends the range where cancellation is proved from
the first extreme into a small neighbourhood of the corner $x=y=1/2$.

\subsection{Provenance}
The determinant-sum estimate was first obtained by a considerably longer
route in an exploratory, AI-generated research notebook.  That route went
through selected-class CRT expansions, a double Heath-Brown decomposition,
and a modulus-averaged reciprocal-energy argument with pair-coprimality
projectors.

The argument above is a streamlined reorganization.  Poisson summation
followed by a single Cauchy--Schwarz over the shorter M\"obius variable
reduces the problem directly to Theorem~\ref{thm:meanvalue}.  This removes
the need for the pair-coprimality projector.  The gcd stratification by
$g'=(a,u)$ in Step~3 replaces a divisor projector used in the notebook.

Sections~\ref{sec:level} and~\ref{sec:LambdaR} are new in this manuscript:
the correlation form, the parity and orientation reduction, the smoothing
argument and the main-term evaluation.  These are among the points that
most need independent verification.

Section~\ref{sec:barrier} was written after a literature check showed that
$9/17$ is not the best exponent available.  It records an AI-driven attempt
to pass $4/7$, made on 15 November 2026 at the investigator's direction,
including the routes that failed.  Nothing in it is used in the proofs of
Theorems~\ref{thm:LambdaR}, \ref{thm:level} and~\ref{thm:main}.

\appendix

\section{The Bettin--Chandee range}\label{app:BC}

For coefficient sequences $\alpha,\beta,\nu$ supported on $[M/2,M]$,
$[N/2,N]$ and $[A_0/2,A_0]$, and $\vartheta\ne0$, let
\[
 \cB(M,N,A_0)=\sum_{\substack{a,m,n\\(m,n)=1}}\alpha_m\beta_n\nu_a\,
 e\Bigl(\vartheta\frac{a\inv m}{n}\Bigr),
\]
where $\inv m$ is an inverse of $m$ modulo $n$.  The bound below is the one
from which Bettin and Chandee deduce their Theorem~1.  It is displayed at
the start of the section ``Completion of the proof of Theorem~1''
of~\cite{BC} (arXiv:1502.00769):
\begin{equation}\label{eq:BC}
\begin{split}
 \cB(M,N,A_0)\ll{}&\|\alpha\|\|\beta\|\|\nu\|(A_0MN)^{\eps}
 \Bigl(1+\frac{|\vartheta|A_0}{MN}\Bigr)^{1/4}\\
 &\times\bigl(A_0^{1/2}M^{1/2}N^{3/8}+A_0^{1/2}N^{7/8}
 +A_0^{1/5}M^{3/5}N^{7/20}+A_0^{7/20}M^{3/10}N^{13/20}\bigr).
\end{split}
\end{equation}

Apply this to the nonzero frequencies in~\eqref{eq:after-poisson}, for
$x\le y$ in the hard square, after the separation of variables of
Section~\ref{sec:proof}.  Take
\[
 m=d,\quad n=e,\quad a=h,\quad M=D,\quad N=E,\quad A_0=H,\quad
 \vartheta=-\kappa/2 .
\]
The coefficient norms are $\|\alpha\|\ll D^{1/2}$, $\|\beta\|\ll E^{1/2}$ and
$\|\nu\|\ll H^{1/2}$.  The Poisson prefactor is $A/E$, and
$|\vartheta|H/(DE)\ll T/(DE)\asymp X^{-1}$.  Since $AD=X$, the frequency
block $|h|\asymp H$ is bounded by
\[
 X^{1+\eps}\bigl(HE^{-1/8}+HE^{3/8}D^{-1/2}
 +H^{7/10}D^{1/10}E^{-3/20}+H^{17/20}D^{-1/5}E^{3/20}\bigr).
\]
Every power of $H$ is positive, so the worst block is $H=TX^{\eps}$.
\begin{itemize}
\item At $H=T$, the four terms are $X^{-\delta}$ for some fixed $\delta>0$
precisely when, respectively,
\[
 8x+7y>7,\qquad 4x+11y>7,\qquad 16x+11y>13,\qquad 13x+20y>16 .
\]
\item In the ordered half-square, the first two conditions follow from the
third.
\item The opposite Poisson orientation adds nothing.  This verification is
routine and was also carried out in the exploratory notebook.
\end{itemize}
This gives the region of Remark~\ref{rem:comparison}; on the diagonal it is
$x>16/33$.  The argument uses no property of $\mu$ beyond $|\mu|\le1$.  It
is therefore what the general-coefficient trilinear estimate gives for
these sums.  The same pair of exponents, $17/33$ from~\eqref{eq:BC} against
$48/95$ from Duke--Friedlander--Iwaniec, appears in the application of
Bettin, Chandee and Radziwi\l\l{} to the twisted second moment of the zeta
function~\cite{BCR}.

\section*{AI usage disclosure}

This manuscript was produced with extensive use of generative AI, and this
disclosure is intended to meet the policies of journals that require one.
AI systems are not authors of this work.  The investigator, Zach Conn,
directed the work and takes responsibility for it.

\begin{itemize}
\item \emph{Origin of the mathematics.}  The determinant-sum estimate, the
key lemmas, and their first proofs were generated by AI research agents
working in an exploratory notebook under the investigator's direction.
[\emph{Investigator to specify the tools.}  The notebook's metadata indicates
OpenAI ChatGPT for the initial handoff reports and an AI coding agent,
apparently OpenAI Codex, for the notebook itself.]
\item \emph{This manuscript.}  The present exposition was drafted on
14--15 November 2026 by Claude Opus 5.5 (Anthropic) running in Claude Code.
This includes:
\begin{itemize}
\item the reorganization of the proof into the direct route described in
Section~\ref{sec:remarks};
\item Sections~\ref{sec:level} and~\ref{sec:LambdaR} and
Appendix~\ref{app:BC};
\item the statements and proofs of all lemmas as written;
\item the literature search and the bibliography.
\end{itemize}
The Bettin--Chandee bound~\eqref{eq:BC} and the exponent $48/95$ of
Coppola--Laporta were checked against the authors' arXiv sources.
Independent AI reviewer instances re-derived the central exponent
computation.  Finite numerical checks supplement the proofs but are not
part of them.  The checks covered the identity~\eqref{eq:pair-phase}, the
gcd expansion of Lemma~\ref{lem:gcd}, the convolution identity in
Lemma~\ref{lem:Sg}, and the convergence $M(R)\to\mathfrak S$.
\item \emph{The barrier section.}  Section~\ref{sec:barrier}, the working
notes and experiments in the repository (\texttt{notes/}), the file
\texttt{Barrier.lean}, and the repositioning of this paper as a
formalization and synthesis were produced on 15 November 2026 by Claude
Fable 5.2 (Anthropic) running in Claude Code, at the investigator's
direction.  This includes the two-sided reduction, the equivalent forms of
the missing estimate, the no-go count, the two numerical experiments, and
the reading of~\cite{BBLR,Polymath8a,FKM3,But,May2,Pas25,GM24,GM25,Lau} from
their arXiv sources.  Every Claim in that section is an AI computation that
no human has checked.
\item \emph{Lean formalization.}  The Lean code of this project
(Section~\ref{sec:lean}) was written by Claude Opus 5.5 instances on 15
November 2026, and \texttt{Barrier.lean} by Claude Fable 5.2 on the same
day; all of it is checked by the Lean kernel.  It imports Lean Pool and
PNT+.  Lean Pool's Kloosterman--Weil development states that generative AI
contributed substantially to it.  Machine checking guarantees that the
formal theorem follows from the formal definitions.  It does not guarantee
that those definitions match the mathematics; that correspondence still
requires human review.
\item \emph{Verification status.}  As of this draft, \textbf{no human
expert has verified the proof}.  The formalization covers
Theorems~\ref{thm:LambdaR}, \ref{thm:level}, \ref{thm:main}
and~\ref{thm:meanvalue} and Corollary~\ref{cor:diagonal}, but not the prose
of this paper.  Not formalized are the comparison in
Appendix~\ref{app:BC}, the comparisons with the literature in
Section~\ref{sec:intro} and Remark~\ref{rem:comparison}, the Claims of
Section~\ref{sec:barrier} (only their arithmetic is checked), and the
correspondence between the Lean definitions and the paper's.
\item \emph{Literature.}  The literature review of 15 November 2026 was
carried out by a Claude Opus 5.5 research agent.  It read the arXiv sources
of~\cite{Walker,Irv14,BCR,PR18,FR22} and checked the bibliographic data
against Crossref and MathSciNet reference matching.  The investigator has not
yet checked it.  All bibliographic details and all statements attributed to
the literature were assembled with AI assistance.  This applies especially
to Lemmas~\ref{lem:weil}, \ref{lem:mobius} and~\ref{lem:BM}, to the
comparisons in Section~\ref{sec:intro} and Remark~\ref{rem:comparison}, and
to Appendix~\ref{app:BC}.  They must be checked against the original sources
before submission or citation.  The transposition of Irving's estimate to
$\mu$ in Remark~\ref{rem:comparison} and the transposition of Conrey's
method in Claim~\ref{claim:conrey} are AI computations that have not been
checked in detail.  The bibliographic data of the references added for
Section~\ref{sec:barrier} were taken from arXiv listings and from memory of
the published versions, and need the same check.
\end{itemize}

\begin{thebibliography}{99}

\bibitem{AB}
M.~C.~Aydemir and M.~Boran,
\emph{Improved averaged distribution of $d_3(n)$ in prime arithmetic
progressions},
arXiv:2601.12601.

\bibitem{BCHB}
R.~Balasubramanian, J.~B.~Conrey and D.~R.~Heath-Brown,
\emph{Asymptotic mean square of the product of the Riemann zeta-function and
a Dirichlet polynomial},
J. Reine Angew. Math. \textbf{357} (1985), 161--181.

\bibitem{BBLR}
S.~Bettin, H.~M.~Bui, X.~Li and M.~Radziwi\l\l,
\emph{A quadratic divisor problem and moments of the Riemann zeta-function},
J. Eur. Math. Soc. \textbf{22} (2020), no.~12, 3953--3980.
arXiv:1609.02539.

\bibitem{BC}
S.~Bettin and V.~Chandee,
\emph{Trilinear forms with Kloosterman fractions},
Adv. Math. \textbf{328} (2018), 1234--1262.  arXiv:1502.00769.

\bibitem{BCR}
S.~Bettin, V.~Chandee and M.~Radziwi\l\l,
\emph{The mean square of the product of the Riemann zeta-function with
Dirichlet polynomials},
J. Reine Angew. Math. \textbf{729} (2017), 51--79.  arXiv:1411.7764.

\bibitem{BM}
T.~F.~Bloom and J.~Maynard,
\emph{A new upper bound for sets with no square differences},
Compos. Math. \textbf{158} (2022), no.~8, 1777--1798.  arXiv:2011.13266.

\bibitem{BFI}
E.~Bombieri, J.~B.~Friedlander and H.~Iwaniec,
\emph{Primes in arithmetic progressions to large moduli},
Acta Math. \textbf{156} (1986), no.~3--4, 203--251.

\bibitem{But}
J.~Buttcane,
\emph{On sums of hyper-Kloosterman sums},
arXiv:2205.14194.

\bibitem{Con89}
J.~B.~Conrey,
\emph{More than two fifths of the zeros of the Riemann zeta function are on
the critical line},
J. Reine Angew. Math. \textbf{399} (1989), 1--26.

\bibitem{CL}
G.~Coppola and M.~Laporta,
\emph{On the correlations, Selberg integral and symmetry of sieve functions
in short intervals, III},
Mosc. J. Comb. Number Theory \textbf{6} (2016), no.~1, 3--24.
arXiv:1003.0302.

\bibitem{DI82}
J.-M.~Deshouillers and H.~Iwaniec,
\emph{Kloosterman sums and Fourier coefficients of cusp forms},
Invent. Math. \textbf{70} (1982/83), no.~2, 219--288.

\bibitem{DFI}
W.~Duke, J.~Friedlander and H.~Iwaniec,
\emph{Bilinear forms with Kloosterman fractions},
Invent. Math. \textbf{128} (1997), no.~1, 23--43.

\bibitem{DFI95}
W.~Duke, J.~B.~Friedlander and H.~Iwaniec,
\emph{Equidistribution of roots of a quadratic congruence to prime moduli},
Ann. of Math. (2) \textbf{141} (1995), no.~2, 423--441.

\bibitem{Est}
T.~Estermann,
\emph{On Kloosterman's sum},
Mathematika \textbf{8} (1961), no.~1, 83--86.

\bibitem{FKM3}
\'E.~Fouvry, E.~Kowalski and Ph.~Michel,
\emph{On the exponent of distribution of the ternary divisor function},
Mathematika \textbf{61} (2015), no.~1, 121--144.
arXiv:1304.3199.

\bibitem{FR22}
\'E.~Fouvry and M.~Radziwi\l\l,
\emph{Level of distribution of unbalanced convolutions},
Ann. Sci. \'Ec. Norm. Sup\'er. (4) \textbf{55} (2022), no.~2, 537--568.
arXiv:1811.08672.

\bibitem{FS11}
\'E.~Fouvry and I.~E.~Shparlinski,
\emph{On a ternary quadratic form over primes},
Acta Arith. \textbf{150} (2011), no.~3, 285--314.

\bibitem{GM24}
L.~Grimmelt and J.~Merikoski,
\emph{Twisted correlations of the divisor function via discrete averages of
$\mathrm{SL}_2(\mathbb R)$ Poincar\'e series},
arXiv:2404.08502.

\bibitem{GM25}
L.~Grimmelt and J.~Merikoski,
\emph{On the greatest prime factor and uniform equidistribution of quadratic
polynomials},
arXiv:2505.00493.

\bibitem{GPY}
D.~A.~Goldston, J.~Pintz and C.~Y.~Y\i ld\i r\i m,
\emph{Primes in tuples I},
Ann. of Math. (2) \textbf{170} (2009), no.~2, 819--862.

\bibitem{GY}
D.~A.~Goldston and C.~Y.~Y\i ld\i r\i m,
\emph{Higher correlations of divisor sums related to primes I: triple
correlations},
Integers \textbf{3} (2003), Paper A5, 66~pp.

\bibitem{Gur}
R.~Guria,
\emph{Prime solutions to a binary additive equation and mixed moments of
character sums},
arXiv:2410.10856.

\bibitem{HB82}
D.~R.~Heath-Brown,
\emph{Prime numbers in short intervals and a generalized Vaughan identity},
Canad. J. Math. \textbf{34} (1982), no.~6, 1365--1377.

\bibitem{HB86}
D.~R.~Heath-Brown,
\emph{The divisor function $d_3(n)$ in arithmetic progressions},
Acta Arith. \textbf{47} (1986), no.~1, 29--56.

\bibitem{Irv14}
A.~J.~Irving,
\emph{Average bounds for Kloosterman sums over primes},
Funct. Approx. Comment. Math. \textbf{51} (2014), no.~2, 221--235.
arXiv:1301.6372.

\bibitem{IK}
H.~Iwaniec and E.~Kowalski,
\emph{Analytic Number Theory},
Amer. Math. Soc. Colloq. Publ. \textbf{53}, Amer. Math. Soc., Providence,
RI, 2004.

\bibitem{Kor18}
M.~A.~Korolev,
\emph{New estimate for a Kloosterman sum with primes for a composite
modulus},
Mat. Sb. \textbf{209} (2018), no.~5, 54--61; English transl.,
Sb. Math. \textbf{209} (2018), no.~5, 652--659.

\bibitem{KC}
M.~A.~Korolev and M.~E.~Changa,
\emph{New estimate for Kloosterman sums with primes},
Mat. Zametki \textbf{108} (2020), no.~1, 94--101; English transl.,
Math. Notes \textbf{108} (2020), no.~1--2, 87--93.

\bibitem{KMS}
E.~Kowalski, Ph.~Michel and W.~Sawin,
Bilinear forms with Kloosterman sums and applications,
\emph{Ann. of Math.} (2) \textbf{186} (2017), 413--500.

\bibitem{Lau}
C.~F.~Lau,
\emph{Smoothed shifted convolutions of generalised divisor functions},
arXiv:2509.07556.

\bibitem{MRT15}
K.~Matom\"aki, M.~Radziwi\l\l{} and T.~Tao,
\emph{An averaged form of Chowla's conjecture},
Algebra Number Theory \textbf{9} (2015), no.~9, 2167--2196.

\bibitem{May2}
J.~Maynard,
\emph{Primes in arithmetic progressions to large moduli II: Well-factorable
estimates},
arXiv:2006.07088.

\bibitem{MV}
H.~L.~Montgomery and R.~C.~Vaughan,
\emph{Multiplicative Number Theory I. Classical Theory},
Cambridge Stud. Adv. Math. \textbf{97}, Cambridge Univ. Press, Cambridge,
2007.

\bibitem{Pas25}
A.~Pascadi,
\emph{On the exponents of distribution of primes and smooth numbers},
arXiv:2505.00653.

\bibitem{Pas25b}
A.~Pascadi,
Non-abelian amplification and bilinear forms with Kloosterman sums,
preprint, arXiv:2511.08445.

\bibitem{Polymath8a}
D.~H.~J.~Polymath,
\emph{New equidistribution estimates of Zhang type},
Algebra Number Theory \textbf{8} (2014), no.~9, 2067--2199.
arXiv:1402.0811.

\bibitem{PR18}
K.~Pratt and N.~Robles,
\emph{Perturbed moments and a longer mollifier for critical zeros of
$\zeta$},
Res. Number Theory \textbf{4} (2018), no.~1, Paper No.~9.
arXiv:1706.04593.

\bibitem{Sha}
P.~Sharma,
\emph{Bilinear sums with $GL(2)$ coefficients and the exponent of
distribution of $d_3$},
Proc. Lond. Math. Soc., to appear.
arXiv:2303.06087.

\bibitem{Tao16}
T.~Tao,
\emph{The logarithmically averaged Chowla and Elliott conjectures for
two-point correlations},
Forum Math. Pi \textbf{4} (2016), e8, 36~pp.

\bibitem{Toth}
\'A.~T\'oth,
\emph{Roots of quadratic congruences},
Internat. Math. Res. Notices \textbf{2000}, no.~14, 719--739.

\bibitem{Walker}
A.~Walker,
\emph{Correlations of sieve weights and distributions of zeros},
Int. J. Number Theory \textbf{19} (2023), no.~1, 139--167.
arXiv:2101.04418.

\bibitem{Watt}
N.~Watt,
\emph{Kloosterman sums and a mean value for Dirichlet polynomials},
J. Number Theory \textbf{53} (1995), no.~1, 179--210.

\bibitem{Weil}
A.~Weil,
\emph{On some exponential sums},
Proc. Nat. Acad. Sci. U.S.A. \textbf{34} (1948), no.~5, 204--207.

\end{thebibliography}

\end{document}
